\section{Constancy of the virtual class along gapped paths}

\subsection{Uniform rigidity of virtual projective classes}

The following finite-dimensional results provide local class constancy when
continuous tensor data and exact unitary covariance are given. They apply
to arbitrary topological parameter spaces, including convergent sequences.
Their source context is the varying-support discussion of
\cite[Appendix~C, lines~2653--2717]{Schuch2011Phases}; no continuous
canonical realization is inferred from a physical spectral gap here.

\begin{definition}[Adjoint action in Hilbert--Schmidt coordinates]
    \label{def:spc_aux_adjoint_coordinates}
    \lean{Matrix.conjugationEuclideanCLM}
    \leanok
    Let $\mathcal H_D=\MN{D}$ have the Hilbert--Schmidt inner product, and
    let $e_D:\mathcal H_D\to\C^{D^2}$ be column vectorization. For
    $V\in\MN{D}$, define the linear operator
    $\beta_V=e_D\circ(X\mapsto VXV^\dagger)\circ e_D^{-1}$.
\end{definition}

\begin{lemma}[Evaluation of the adjoint action]
    \label{thm:spc_aux_adjoint_evaluation}
    \lean{Matrix.conjugationEuclideanCLM_apply}
    \leanok
    \uses{def:spc_aux_adjoint_coordinates}
    For $V,X\in\MN{D}$, one has
    $\beta_V(e_DX)=e_D(VXV^\dagger)$.
\end{lemma}

\begin{proof}\leanok
    \uses{def:spc_aux_adjoint_coordinates}
    Substitute the definitions of vectorization and the adjoint action.
\end{proof}

\begin{definition}[Euclidean coefficient operator]
    \label{def:spc_aux_coefficient_operator}
    \lean{MPSTensor.coefficientEuclideanMap}
    \leanok
    For a tensor $C=\{C^i\}_{i=0}^{d-1}$ with bond dimension $D$, let
    $\Phi_C:\C^d\to\C^{D^2}$ be the operator with matrix entries
    $(\Phi_C)_{(a,b),i}=C^i_{b,a}$.
\end{definition}

\begin{lemma}[Evaluation of the coefficient operator]
    \label{thm:spc_aux_coefficient_evaluation}
    \lean{MPSTensor.coefficientEuclideanMap_apply}
    \leanok
    \uses{def:spc_aux_coefficient_operator}
    For $x\in\C^d$, one has
    $\Phi_Cx=e_D(\sum_i x_i C^i)$.
\end{lemma}

\begin{proof}\leanok
    \uses{def:spc_aux_coefficient_operator}
    Each coordinate is $\sum_i C^i_{b,a}x_i$.
\end{proof}

\begin{lemma}[Continuity of the coefficient operator]
    \label{thm:spc_aux_coefficient_continuity}
    \lean{MPSTensor.continuous_coefficientEuclideanMap}
    \leanok
    \uses{def:spc_aux_coefficient_operator}
    For fixed $d,D$, the map $C\mapsto\Phi_C$ is continuous in the
    operator norm.
\end{lemma}

\begin{proof}\leanok
    \uses{def:spc_aux_coefficient_operator}
    The map is linear between finite-dimensional complex vector spaces.
\end{proof}

\begin{lemma}[Surjectivity of the injective coefficient operator]
    \label{thm:spc_aux_coefficient_surjective}
    \lean{MPSTensor.coefficientEuclideanMap_surjective_of_isInjective}
    \leanok
    \uses{def:spc_aux_coefficient_operator}
    If $\spn_{\C}\{C^i\}=\MN{D}$, then $\Phi_C$ is surjective.
\end{lemma}

\begin{proof}\leanok
    \uses{thm:spc_aux_coefficient_evaluation}
    Every matrix is a linear combination of the letters. Apply vectorization
    to that combination.
\end{proof}

\begin{definition}[Gram right inverse of the coefficient operator]
    \label{def:spc_aux_coefficient_inverse}
    \lean{MPSTensor.coefficientRightInverse}
    \leanok
    \uses{def:spc_aux_coefficient_operator}
    Define $\Psi_C=\Phi_C^\dagger(\Phi_C\Phi_C^\dagger)^{-1}$.
    Here the inverse is assigned the value zero when the operator is not
    invertible.
\end{definition}

\begin{lemma}[Right-inverse identity at an injective tensor]
    \label{thm:spc_aux_coefficient_inverse_identity}
    \lean{MPSTensor.coefficientEuclideanMap_comp_rightInverse_of_isInjective}
    \leanok
    \uses{def:spc_aux_coefficient_inverse}
    If $\spn_{\C}\{C^i\}=\MN{D}$, then $\Phi_C\Psi_C=I$.
\end{lemma}

\begin{proof}\leanok
    \uses{thm:spc_aux_coefficient_surjective, def:spc_aux_coefficient_inverse}
    Surjectivity makes $\Phi_C^\dagger$ injective. Thus
    $\Phi_C\Phi_C^\dagger$ is invertible, and $\Phi_C\Psi_C=I$ by
    inverse cancellation.
\end{proof}

\begin{lemma}[Continuity of the Gram right inverse]
    \label{thm:spc_aux_coefficient_inverse_continuity}
    \lean{MPSTensor.continuousAt_coefficientRightInverse_of_isInjective}
    \leanok
    \uses{def:spc_aux_coefficient_inverse}
    If $\spn_{\C}\{C^i\}=\MN{D}$, the map $B\mapsto\Psi_B$ is
    continuous at $B=C$.
\end{lemma}

\begin{proof}\leanok
    \uses{thm:spc_aux_coefficient_continuity, thm:spc_aux_coefficient_surjective, def:spc_aux_coefficient_inverse}
    The Gram operator $\Phi_B\Phi_B^\dagger$ is continuous in $B$ and
    invertible at $C$. Inversion is continuous at invertible operators,
    so its product with $\Phi_B^\dagger$ is continuous at $C$.
\end{proof}

\begin{lemma}[Coefficient form of exact covariance]
    \label{thm:spc_aux_coefficient_covariance}
    \lean{MPSTensor.coefficientEuclideanMap_covariance}
    \leanok
    \uses{def:spc_aux_adjoint_coordinates, def:spc_aux_coefficient_operator}
    Let $U\in\MN{d}$ and $V\in\MN{D}$ satisfy
    $\sum_j U_{ij}C^j=VC^iV^\dagger$ for every $i$. Then
    $\Phi_C U^{\mathsf T}=\beta_V\Phi_C$.
    Neither matrix is required to be unitary or invertible.
\end{lemma}

\begin{proof}\leanok
    \uses{thm:spc_aux_coefficient_evaluation, thm:spc_aux_adjoint_evaluation}
    For $x\in\C^d$, expand the two sums to obtain
    \begin{align}
        \sum_j (U^{\mathsf T}x)_j C^j
         & =\sum_i x_i\sum_j U_{ij}C^j
        =V\Bigl(\sum_i x_iC^i\Bigr)V^\dagger.\notag
    \end{align}
    Apply vectorization.
\end{proof}

\begin{lemma}[Adjoint action determined by the physical action]
    \label{thm:spc_aux_adjoint_coefficient_formula}
    \lean{MPSTensor.conjugationEuclideanCLM_eq_coefficients_of_isInjective}
    \leanok
    \uses{def:spc_aux_adjoint_coordinates, def:spc_aux_coefficient_inverse}
    Suppose $C$ is one-site injective and
    $\sum_j U_{ij}C^j=VC^iV^\dagger$ for every $i$, with
    $U\in\MN{d}$ and $V\in\MN{D}$ arbitrary. Then
    $\beta_V=\Phi_C U^{\mathsf T}\Psi_C$.
\end{lemma}

\begin{proof}\leanok
    \uses{thm:spc_aux_coefficient_covariance, thm:spc_aux_coefficient_inverse_identity}
    Multiply the coefficient covariance identity on the right by
    $\Psi_C$ and use $\Phi_C\Psi_C=I$.
\end{proof}

\begin{lemma}[Uniform adjoint estimate near an injective tensor]
    \label{thm:spc_aux_uniform_adjoint_estimate}
    \lean{MPSTensor.eventually_uniform_conjugationEuclideanCLM_of_covariance}
    \leanok
    \uses{def:spc_aux_adjoint_coordinates, def:spc_aux_coefficient_operator}
    Let $C_0$ be one-site injective, let $\mathcal I$ be an arbitrary set,
    and let $U_g\in\mathrm U(d)$ and $V_g^0\in\MN{D}$ satisfy
    $\sum_j(U_g)_{ij}C_0^j=V_g^0C_0^i(V_g^0)^\dagger$ for all $g,i$.
    For every $\varepsilon>0$ there is a neighborhood of $C_0$ such that,
    for each tensor $C$ in that neighborhood, every $g\in\mathcal I$ and
    every $V\in\MN{D}$ satisfying
    $\sum_j(U_g)_{ij}C^j=VC^iV^\dagger$ obey
    $\|\beta_V-\beta_{V_g^0}\|<\varepsilon$.
    No continuity or unitarity of the virtual matrices is assumed.
\end{lemma}

\begin{proof}\leanok
    \uses{thm:spc_aux_coefficient_continuity, thm:spc_aux_coefficient_inverse_continuity, thm:spc_aux_coefficient_surjective, thm:spc_aux_coefficient_covariance, thm:spc_aux_adjoint_coefficient_formula}
    Nearby Gram operators remain invertible. Exact covariance gives
    $\beta_V=\Phi_CU_g^{\mathsf T}\Psi_C$ there. Since
    $\|U_g^{\mathsf T}\|\leq1$, the difference has the bound
    \begin{align}
        \|\Phi_CU_g^{\mathsf T}\Psi_C-
        \Phi_{C_0}U_g^{\mathsf T}\Psi_{C_0}\|
         & \leq \|\Phi_C-\Phi_{C_0}\|\,\|\Psi_C\|
        +\|\Phi_{C_0}\|\,\|\Psi_C-\Psi_{C_0}\|.\notag
    \end{align}
    Continuity of the coefficient operator and its right inverse makes
    this bound uniformly smaller than $\varepsilon$.
\end{proof}

\begin{lemma}[Cancellation of a unitary projective factor]
    \label{thm:spc_aux_unitary_factor_cancel}
    \lean{TNLean.Algebra.ProjectiveRepresentation.factor_mul_star_eq_one_of_unitary}
    \leanok
    \uses{def:projective_rep}
    Let $\rho:G\to\mathrm{GL}(D,\C)$ be a projective representation
    with factor $\omega$, where $D>0$. If every $\rho(g)$ is unitary,
    then $\omega(g,h)\overline{\omega(g,h)}=1$ for every $g,h$.
\end{lemma}

\begin{proof}\leanok
    \uses{thm:spt_unitary_factor_alignment}
    Unitary projective factors have modulus one. Hence
    $\omega(g,h)\overline{\omega(g,h)}=|\omega(g,h)|^2=1$.
\end{proof}

\begin{definition}[Projective action on the intertwiner space]
    \label{def:spc_aux_projective_hom_action}
    \lean{TNLean.Algebra.ProjectiveRepresentation.homMatrix}
    \leanok
    \uses{def:projective_rep}
    For projective representations $\rho_0:G\to\mathrm{GL}(D_0,\C)$
    and $\rho_1:G\to\mathrm{GL}(D_1,\C)$, set
    $S_g=\overline{\rho_0(g)}\otimes\rho_1(g)$ on
    $\C^{D_0}\otimes\C^{D_1}$. Under column vectorization this is
    $X\mapsto\rho_1(g)X\rho_0(g)^\dagger$ on rectangular matrices.
\end{definition}

\begin{lemma}[Unitarity on the intertwiner space]
    \label{thm:spc_aux_projective_hom_unitary}
    \lean{TNLean.Algebra.ProjectiveRepresentation.homMatrix_mem_unitaryGroup}
    \leanok
    \uses{def:spc_aux_projective_hom_action}
    If all matrices of $\rho_0$ and $\rho_1$ are unitary, then every
    $S_g$ is unitary. The dimensions $D_0,D_1$ may be zero.
\end{lemma}

\begin{proof}\leanok
    \uses{def:spc_aux_projective_hom_action}
    Entrywise conjugation preserves unitarity, as does the tensor product.
\end{proof}

\begin{lemma}[Multiplication on the intertwiner space]
    \label{thm:spc_aux_projective_hom_product}
    \lean{TNLean.Algebra.ProjectiveRepresentation.homMatrix_mul}
    \leanok
    \uses{def:spc_aux_projective_hom_action}
    If $\rho_a$ has factor $\omega_a$, then
    $S_gS_h=\overline{\omega_0(g,h)}\omega_1(g,h)S_{gh}$ for all $g,h$.
    Unitarity is not required.
\end{lemma}

\begin{proof}\leanok
    \uses{def:spc_aux_projective_hom_action}
    Use multiplicativity of entrywise conjugation and
    $(A\otimes B)(C\otimes D)=AC\otimes BD$, then substitute the two
    projective multiplication laws.
\end{proof}

\begin{lemma}[Exact composition of conjugations on intertwiners]
    \label{thm:spc_aux_projective_hom_conjugation}
    \lean{TNLean.Algebra.ProjectiveRepresentation.homMatrix_conjugation_mul}
    \leanok
    \uses{def:spc_aux_projective_hom_action}
    Suppose $D_0,D_1>0$ and both projective representations are unitary.
    For every $g,h\in G$ and matrix $M$ on the intertwiner space,
    $S_g(S_hMS_h^\dagger)S_g^\dagger=S_{gh}MS_{gh}^\dagger$.
    The two factor systems need not agree.
\end{lemma}

\begin{proof}\leanok
    \uses{thm:spc_aux_projective_hom_product, thm:spc_aux_unitary_factor_cancel}
    Substitute $S_gS_h=\lambda_{g,h}S_{gh}$, where
    $\lambda_{g,h}=\overline{\omega_0(g,h)}\omega_1(g,h)$.
    The scalar cancels because $\lambda_{g,h}\overline{\lambda_{g,h}}=1$.
\end{proof}

\begin{lemma}[Uniform rigidity of a unitary projective class]
    \label{thm:spc_aux_uniform_projective_rigidity}
    \lean{TNLean.Algebra.ProjectiveRepresentation.exists_uniform_adjoint_rigidity_radius}
    \leanok
    \uses{def:spc_aux_adjoint_coordinates, def:projective_rep, def:cohomologous_to}
    For every $D>0$ there is $\delta_D>0$ such that the following holds
    for any group $G$. Let $\rho_0,\rho_1$ be unitary projective
    representations on $\C^D$, with factors $\omega_0,\omega_1$.
    If $\|\beta_{\rho_1(g)}-\beta_{\rho_0(g)}\|<\delta_D$ for every $g$,
    then $\omega_1$ is cohomologous to $\omega_0$.
    No topology on $G$ or continuity of either representation is required.
\end{lemma}

\begin{proof}\leanok
    \uses{def:spc_aux_projective_hom_action, thm:spc_aux_projective_hom_unitary, thm:spc_aux_projective_hom_conjugation, thm:spt_local_spectral_support, thm:spt_virtual_class_gauge_preparation}
    Let $J_0$ span the normalized vectorization of the identity matrix
    and let $P_0=J_0J_0^\dagger$. The linear rearrangement of matrix entries
    which takes an adjoint operator to its rank-one Choi matrix shows that
    uniformly close adjoint actions give uniformly close conjugates
    $S_gP_0S_g^\dagger$. The operators $S_g$ are unitary and their
    conjugations compose exactly. In the closed convex hull of the orbit
    of $P_0$, the unique least-norm element $R$ is therefore fixed by every
    conjugation. It is Hermitian and remains close to $P_0$.

    Choose the radius so that the spectral line of $R$ continuing the
    range of $P_0$ has a unit vector $J$ whose decoded matrix $W$ is
    invertible. Every $S_g$ preserves this line, giving
    $\rho_1(g)W=c(g)W\rho_0(g)$. Invertibility implies $c(g)\neq0$.
    Comparing multiplication at $g,h$ gives
    $\omega_1(g,h)=c(g)c(h)c(gh)^{-1}\omega_0(g,h)$, as required.
\end{proof}

\begin{lemma}[Local constancy from continuous tensors and exact covariance]
    \label{thm:spc_aux_local_tensor_class}
    \lean{MPSTensor.eventually_cohomologousTo_of_continuousAt_exact_unitary_covariance}
    \leanok
    \uses{def:projective_rep, def:cohomologous_to}
    Let $T$ be a topological space, let $t_0\in T$, and let
    $C:T\to(\MN{D})^d$ be continuous at $t_0$, with $D>0$ and $C(t_0)$
    one-site injective. Fix a group $G$ and matrices $U_g\in\mathrm U(d)$.
    For each $t$, let $\rho_t$ be a unitary projective representation with
    factor $\omega_t$, and suppose
    \begin{align}
        \sum_j(U_g)_{ij}C^j(t)
         & =\rho_t(g^{-1})C^i(t)\rho_t(g^{-1})^\dagger.\notag
    \end{align}
    Then $\omega_t$ is cohomologous to $\omega_{t_0}$ throughout some
    neighborhood of $t_0$. No continuity of $\rho_t$ or $\omega_t$,
    injectivity away from $t_0$, or connectedness of $T$ is assumed.
\end{lemma}

\begin{proof}\leanok
    \uses{thm:spc_aux_uniform_adjoint_estimate, thm:spc_aux_uniform_projective_rigidity}
    Use the uniform adjoint estimate at $C(t_0)$ with tolerance
    $\delta_D$, and pull its tensor neighborhood back along $C$.
    Reindex by $g^{-1}$ to obtain
    $\|\beta_{\rho_t(g)}-\beta_{\rho_{t_0}(g)}\|<\delta_D$ for every $g$.
    Apply uniform projective rigidity at each parameter in that neighborhood.
\end{proof}


\begin{lemma}[Local class constancy from stationary support data]
    \label{thm:spc_aux_topological_stationary_support_class}
    \lean{MPSTensor.eventually_cohomologousTo_of_continuousAt_unital_supported_family}
    \leanok
    \uses{def:projective_rep, def:cohomologous_to, def:transfer_map}
    Let $T$ be a topological space and $t_0\in T$. Let $B(t)$ be a unital
    tensor family of fixed physical dimension $d$ and ambient bond
    dimension $k$, continuous at $t_0$. Suppose that for every $t$ there
    are a positive semidefinite matrix $\sigma_t$, a dimension $D_t$,
    and an isometry $K_t:\C^{D_t}\to\C^k$ satisfying
    \begin{align}
        \tr\sigma_t                   & =1,                                   &
        \mathcal E_{B(t)}^*(\sigma_t) & =\sigma_t,                            &
        K_tK_t^\dagger                & =\operatorname{supp}(\sigma_t).\notag
    \end{align}
    Put $A_t^i=K_t^\dagger B(t)^iK_t$. Let $U_g$ be a fixed unitary
    physical representation, and let $V_{t,g}$ be unitary projective
    representations on $\C^{D_t}$, with factor systems $\omega_t$, such that
    \begin{align}
        \sum_j(U_g)_{ij}A_t^j
         & =V_{t,g^{-1}}A_t^iV_{t,g^{-1}}^\dagger.\notag
    \end{align}
    If $A_{t_0}$ is one-site injective and
    $\dim\ker(I-\mathcal E_{B(t_0)}^*)=1$, then
    $[\omega_t]=[\omega_{t_0}]$ for all $t$ in some neighborhood of $t_0$.
    The stationary matrices, full support frames, and virtual
    representations are not required to be continuous. The ambient
    tensors and their supported covariance are assumed;
    this statement does not construct them from a physical gap.
\end{lemma}

\begin{proof}\leanok
    \uses{thm:spt_local_unique_stationary_class}
    Refine the topology of $T$ by the topology induced by $B$.
    Then $B$ is continuous everywhere, and continuity at $t_0$ implies
    that the neighborhood filter at $t_0$ is unchanged. It therefore
    suffices to prove the assertion in this finer topology.

    The normalized stationary equation is invertible near $t_0$.
    Uniqueness identifies the given matrices $\sigma_t$ with the
    continuous normalized stationary branch there. Choose positive
    spectral cutoffs retaining exactly the support of $\sigma_{t_0}$.
    The resulting spectral projections have a continuous isometric
    frame $J_t:\C^{D_{t_0}}\to\C^k$, with $J_{t_0}=K_{t_0}$.

    Uniqueness of the normalized stationary matrix on the full support
    implies that every $V_{t,g}$ commutes with
    $K_t^\dagger\sigma_tK_t$. Extending each such matrix by the identity
    on the orthogonal complement shows that it preserves the selected
    ambient spectral corner. This extension is used individually;
    no projective representation on the ambient complement is needed.
    Since the corner is contained in the full support, $K_t^\dagger J_t$
    is an invariant isometry. Restriction along this isometry retains
    the factor system $\omega_t$ and gives exact unitary covariance for
    $C_t^i=J_t^\dagger B(t)^iJ_t$. The tensor letters need not preserve
    the corner for this covariance identity.

    The family $C_t$ is continuous near $t_0$, and
    $C_{t_0}=A_{t_0}$ is injective. Local rigidity under exact unitary
    covariance gives the required equality of cohomology classes.
    Transfer back through the unchanged neighborhood filter at $t_0$.
\end{proof}


\subsection{Compact limits and transient periodic components}

The following auxiliary results apply when a uniform bound below one on
nonunit transfer eigenvalues is given. They permit a limiting stationary
density to lose rank, and control the periodic contribution of the
complementary corner. The derivation of this transfer bound from a physical
Hamiltonian gap remains separate
\cite{gap:spc11_spt_interpolation_upper_range}.


\begin{lemma}[Compactness of unital tensors]
    \label{thm:spc_aux_compact_unital_tensors}
    \lean{MPSTensor.isCompact_setOf_isUnital}
    \leanok
    \uses{def:mps_tensor}
    For fixed $d,D\geq 0$, the set of tensors $A=(A^i)_{i=1}^d$ satisfying
    $\sum_i A^i(A^i)^\dagger=I_D$ is compact in the topology of their
    matrix entries. No injectivity or stationary-density assumption is needed.
\end{lemma}

\begin{proof}\leanok
    The normalization condition is closed. Its diagonal entries give
    $\sum_{i,c}|A^i_{ac}|^2=1$, so every matrix entry has modulus at most
    one. The set is therefore a closed subset of a finite product of
    compact complex unit discs.
\end{proof}

\begin{lemma}[A common stationary subsequence]
    \label{thm:spc_aux_unital_stationary_subsequence}
    \lean{MPSTensor.exists_subsequence_unital_stationary_limit}
    \leanok
    \uses{def:transfer_map}
    Let $A_n$ be unital tensors of fixed physical dimension $d$ and bond
    dimension $D$. Suppose that $\sigma_n\geq 0$, $\tr\sigma_n=1$, and
    $\mathcal E_{A_n}^*(\sigma_n)=\sigma_n$, where
    $\mathcal E_A^*(X)=\sum_i(A^i)^\dagger X A^i$.
    There exist a strictly increasing sequence of indices $n_k$, a unital
    tensor $B$, and a density matrix $\tau$ such that
    $A_{n_k}\to B$, $\sigma_{n_k}\to\tau$, and
    $\mathcal E_B^*(\tau)=\tau$. The rank of $\tau$ may be smaller than
    the ranks of the matrices $\sigma_n$.
\end{lemma}

\begin{proof}\leanok
    \uses{thm:spc_aux_compact_unital_tensors}
    Take a convergent subsequence in the product of the compact set of
    unital tensors and the compact set of density matrices. Matrix
    multiplication and adjoints are continuous, so stationarity passes
    to the limit.
\end{proof}

\begin{lemma}[Convergence on a fixed invariant quotient]
    \label{thm:spc_aux_quotient_operator_limit}
    \lean{Submodule.tendsto_toMatrix_mapQ_of_tendsto}
    \leanok
    Let $V$ be a finite-dimensional complex normed space, let
    $f_n\to f_\infty$ in operator norm, and let $W\subseteq V$ be invariant
    under every $f_n$ and under $f_\infty$. In any fixed basis of $V/W$,
    the matrices of the induced quotient operators converge to the
    matrix of the quotient operator induced by $f_\infty$.
\end{lemma}

\begin{proof}\leanok
    Choose a representative in $V$ of each quotient basis vector. Each
    matrix entry is a fixed linear functional applied to the image of
    one representative under the operator. These functionals are
    continuous in finite dimension, so their values converge.
\end{proof}

\begin{lemma}[Stability of a fixed line under spectral contraction]
    \label{thm:spc_aux_fixed_line_limit}
    \lean{ContinuousLinearMap.eigenspace_one_limit_eq_span_of_uniform_spectrum_bound}
    \leanok
    Let $f_n\to f_\infty$ in operator norm on a finite-dimensional complex
    normed space $V$. Suppose $v\neq 0$, $f_n v=v$ for every $n$, and the
    eigenvalue $1$ has algebraic multiplicity one for every $f_n$.
    If there exists $q<1$ such that every eigenvalue $\lambda\neq 1$ of
    every $f_n$ satisfies $|\lambda|\leq q$, then
    $\ker(f_\infty-I)=\spn_{\C}\{v\}$.
    The common spectral bound is an assumption of this statement.
\end{lemma}

\begin{proof}\leanok
    \uses{thm:spc_aux_quotient_operator_limit}
    Continuity gives $f_\infty v=v$. Put $W=\spn_{\C}\{v\}$ and let
    $Q_n,Q_\infty$ denote the induced operators on $V/W$, in one fixed
    basis. Algebraic simplicity removes exactly the eigenvalue $1$
    when passing from $f_n$ to $Q_n$. Thus all eigenvalues of $Q_n$ have
    modulus at most $q$. If $m=\dim(V/W)$, the characteristic polynomial
    gives
    \begin{align}
        |\det(I-Q_n)|
         & =\prod_{j=1}^m|1-\lambda_j(Q_n)|
        \geq(1-q)^m>0.\notag
    \end{align}
    Quotient coordinates and determinants converge, so $I-Q_\infty$ is
    invertible. Every fixed vector of $f_\infty$ therefore has zero
    image in $V/W$, proving the assertion. No diagonalizability is used.
\end{proof}

\begin{lemma}[A nontrivial spectral bound in the limit]
    \label{thm:spc_aux_nontrivial_spectrum_limit}
    \lean{ContinuousLinearMap.spectrum_limit_norm_le_of_uniform_nontrivial_bound}
    \leanok
    Let $f_n\to f_\infty$ in operator norm on a finite-dimensional complex
    normed space. Suppose $v\neq 0$, $f_n v=v$ for every $n$, and the
    eigenvalue $1$ has algebraic multiplicity one for every $f_n$.
    Let $q\in\mathbb R$ satisfy $|\lambda|\leq q$ for every eigenvalue
    $\lambda\neq 1$ of every $f_n$. Then every eigenvalue
    $\lambda\neq 1$ of $f_\infty$ also satisfies $|\lambda|\leq q$.
    Neither $q<1$ nor algebraic simplicity of $1$ at the limit is required.
\end{lemma}

\begin{proof}\leanok
    \uses{thm:spc_aux_quotient_operator_limit}
    Pass to the fixed quotient by $\spn_{\C}\{v\}$, as in the preceding
    fixed-line argument. Its operators $Q_n$ have all eigenvalues in
    the closed disc of radius $q$. For $|z|>q$, their characteristic
    polynomials give
    \begin{align}
        |\det(zI-Q_n)|\geq(|z|-q)^{\dim(V/\spn_{\C}\{v\})}>0.\notag
    \end{align}
    This lower bound survives the limit, so $z$ is absent from the
    quotient spectrum of $f_\infty$. An eigenvector of $f_\infty$ with
    eigenvalue different from $1$ has nonzero image in this quotient.
    Its eigenvalue therefore lies in the closed disc as asserted.
\end{proof}

\begin{lemma}[Equality of nonzero limiting periodic lines]
    \label{thm:spc_aux_periodic_line_limit}
    \lean{MPSTensor.mpv_span_eq_of_tendsto}
    \leanok
    \uses{def:mpv}
    Fix a physical dimension $d$, bond dimensions $D,E$, and a length
    $N\geq 0$. Let $A_x\to A$ and $B_x\to B$ along a nontrivial filter,
    with bond dimensions $D$ and $E$, respectively. Suppose that,
    eventually along this filter,
    $\spn_{\C}\{\ket{V^{(N)}(A_x)}\}
        =\spn_{\C}\{\ket{V^{(N)}(B_x)}\}$.
    If both limiting periodic vectors are nonzero, then
    $\spn_{\C}\{\ket{V^{(N)}(A)}\}
        =\spn_{\C}\{\ket{V^{(N)}(B)}\}$.
\end{lemma}

\begin{proof}\leanok
    \uses{def:eval_word}
    Periodic-vector coordinates are traces of finite products, hence
    are continuous in the tensor entries. Equality of the two lines
    forces all two-by-two minors of the pair of coordinate vectors to
    vanish. These identities pass to the limit. Choosing a nonzero
    coordinate of the first limiting vector shows that the second is
    a scalar multiple of it; nonvanishing of the second vector makes
    that scalar nonzero.
\end{proof}


\begin{lemma}[Deviation of the periodic MPS norm]
    \label{thm:spt_compact_periodic_norm_deviation}
    \lean{MPSTensor.abs_norm_mpvState_sq_sub_one_le_of_transfer_spectrum}
    \leanok
    Let $A$ have physical dimension $d$ and bond dimension $D>0$.
    Suppose that $\E_A(\Id)=\Id$ and $\dim\ker(\E_A-\Id)=1$.
    Let $q\in\R$ satisfy $|\lambda|\leq q$ for every transfer eigenvalue
    $\lambda\ne1$. Then every integer $N\geq0$ satisfies
    \begin{align}
        \bigl|\|\ket{V^{(N)}(A)}\|^2-1\bigr|
         & \leq (D^2-1)q^N.
        \label{eq:spt_compact_norm_deviation}
    \end{align}
\end{lemma}
\begin{proof}\leanok
    \uses{thm:overlap_trace}
    The peripheral Jordan-block theorem for positive unital maps
    \cite[Proposition~6.2]{Wolf2012Quantum} makes the characteristic root
    $1$ algebraically simple. Set $W=\MN{D}/\C\Id$ and let
    $\overline{\E}_A$ be the induced endomorphism. Its spectrum is contained
    in $\operatorname{spec}(\E_A)\setminus\{1\}$, and
    $\operatorname{Tr}(\E_A^N)=1+\operatorname{Tr}(\overline{\E}_A^N)$.
    Spectral mapping and the characteristic-polynomial trace identity give
    $|\operatorname{Tr}(\overline{\E}_A^N)|\leq(\dim W)q^N$,
    with eigenvalues counted according to algebraic multiplicity.
    The periodic overlap identity gives
    $\|\ket{V^{(N)}(A)}\|^2=\operatorname{Tr}(\E_A^N)$.
\end{proof}

\begin{lemma}[Lower bound for the periodic MPS norm]
    \label{thm:spt_compact_periodic_norm_lower}
    \lean{MPSTensor.one_sub_le_norm_mpvState_sq_of_transfer_spectrum}
    \leanok
    Under the hypotheses of
    Lemma~\ref{thm:spt_compact_periodic_norm_deviation}, every integer
    $N\geq0$ satisfies
    $1-(D^2-1)q^N\leq\|\ket{V^{(N)}(A)}\|^2$.
\end{lemma}
\begin{proof}\leanok
    \uses{thm:spt_compact_periodic_norm_deviation}
    Take the lower half of the absolute-value bound
    (\ref{eq:spt_compact_norm_deviation}).
\end{proof}

\begin{lemma}[A uniform lower bound on sufficiently long rings]
    \label{thm:spt_compact_uniform_long_norm}
    \lean{MPSTensor.exists_uniform_norm_mpvState_sq_lowerBound_of_transfer_spectrum}
    \leanok
    Let $\{A_t:t\in T\}$ be an arbitrary family of tensors with common
    dimensions $d$ and $D>0$. Suppose that every $\E_{A_t}$ is unital,
    has a one-dimensional fixed space, and has all eigenvalues other than
    $1$ of modulus at most a common $q$ with $0\leq q<1$.
    There exists an integer $L>0$ such that
    $\|\ket{V^{(N)}(A_t)}\|^2\geq\tfrac12$ for every $t\in T$ and
    every integer $N\geq L$.
\end{lemma}
\begin{proof}\leanok
    \uses{thm:spt_compact_periodic_norm_lower}
    Choose $L>0$ such that $(D^2-1)q^N<\tfrac12$ whenever $N\geq L$,
    and apply the preceding lower bound to each $A_t$.
\end{proof}

\begin{lemma}[Uniform nonvanishing after physical blocking]
    \label{thm:spt_compact_uniform_blocking_nonzero}
    \lean{MPSTensor.exists_uniform_blocking_mpvState_ne_zero_of_transfer_spectrum}
    \leanok
    Under the hypotheses of Lemma~\ref{thm:spt_compact_uniform_long_norm},
    there exists an integer $L>0$ such that every $t\in T$ and every
    integer $N>0$ satisfy
    $\|\ket{V^{(N)}(A_t^{[L]})}\|^2\geq\tfrac12$ and
    $\ket{V^{(N)}(A_t^{[L]})}\ne0$.
\end{lemma}
\begin{proof}\leanok
    \uses{thm:spt_compact_uniform_long_norm}
    Physical blocking gives $\E_{A_t^{[L]}}=\E_{A_t}^L$, hence the
    transfer-trace identity identifies the squared blocked norm at length
    $N$ with the original squared norm at length $LN\geq L$.
    Its positive lower bound implies nonvanishing.
\end{proof}

\begin{lemma}[Transfer spectrum of a complementary diagonal corner]
    \label{thm:spt_compact_transient_corner_spectrum}
    \lean{MPSTensor.spectrum_transferMap_transientCorner_norm_le}
    \leanok
    Let $A$ have dimensions $d$ and $D>0$, with $\E_A$ unital and
    $\dim\ker(\E_A-\Id)=1$. Let $P$ be a nonzero orthogonal projection
    such that $P A^i(\Id-P)=0$ for each physical index $i$.
    Put $R^i=(\Id-P)A^i(\Id-P)$. Suppose that $q\geq0$ and that
    every eigenvalue of $\E_A$ other than $1$ has modulus at most $q$.
    Then every eigenvalue of $\E_R$ has modulus at most $q$.
\end{lemma}
\begin{proof}\leanok
    Set $Q=\Id-P$. If $\E_R(X)=\lambda X$ with $\lambda\ne0$,
    then $QXQ=X$, since the range of $\E_R$ lies in the $Q$-corner.
    The triangular condition gives $\E_A(X)=\E_R(X)$, so $\lambda$
    belongs to the ambient transfer spectrum. If $\lambda=1$, then
    $X=c\Id$ because the ambient fixed space is one-dimensional.
    The support identity forces $cQ=c\Id$, contradicting $P\ne0$
    and $X\ne0$. Thus $\lambda\ne1$; the zero eigenvalue obeys the
    bound because $q\geq0$. The corner decomposition is the one used in
    \cite[proof of the TI canonical-form theorem]{PerezGarcia2007Matrix}.
\end{proof}

\begin{lemma}[Decay of the complementary periodic contribution]
    \label{thm:spt_compact_transient_mpv_decay}
    \lean{MPSTensor.norm_mpvState_transientCorner_tendsto_zero_of_transfer_spectrum}
    \leanok
    Under the hypotheses of
    Lemma~\ref{thm:spt_compact_transient_corner_spectrum}, suppose also
    that $q<1$. Then $\|\ket{V^{(N)}(R)}\|\longrightarrow0$ as
    $N\longrightarrow\infty$.
\end{lemma}
\begin{proof}\leanok
    \uses{thm:spt_compact_transient_corner_spectrum}
    The spectral bound gives $r(\E_R)<1$, so the transfer-trace overlap
    decay theorem gives $O_{RR}(N)\to0$. Since
    $O_{RR}(N)=\|\ket{V^{(N)}(R)}\|^2$, continuity of the square root
    gives the result.
\end{proof}

\begin{lemma}[Vanishing overlap of a periodic remainder]
    \label{thm:spc_aux_overlap_remainder_zero}
    \lean{MPSTensor.mpvOverlap_sub_tendsto_zero_of_eventual_decomposition}
    \leanok
    \uses{def:mpv_state, def:mpv_overlap}
    Let $A,B,C,R$ be tensors with the same physical dimension and arbitrary
    finite bond dimensions. Suppose there is $M\in\R$ such that
    $\|\ket{V^{(N)}(A)}\|\leq M$ for all sufficiently large $N$, and suppose
    $\ket{V^{(N)}(B)}=\ket{V^{(N)}(C)}+\ket{V^{(N)}(R)}$ for all sufficiently large $N$.
    If $\|\ket{V^{(N)}(R)}\|\to0$, then
    $O_{AB}(N)-O_{AC}(N)\to0$.
    This is an auxiliary asymptotic statement for the varying-support
    argument of \cite[Appendix~C, lines~2653--2717]{Schuch2011Phases}.
\end{lemma}

\begin{proof}\leanok
    \uses{lem:overlap_eq_inner, def:mpv_inner, def:mpv_state}
    The decomposition and conjugate-linearity of the second argument of
    the overlap give $O_{AB}(N)-O_{AC}(N)=O_{AR}(N)$ whenever the
    decomposition holds. Cauchy--Schwarz therefore yields
    \begin{align}
        |O_{AB}(N)-O_{AC}(N)|
         & \leq \|\ket{V^{(N)}(A)}\|\,\|\ket{V^{(N)}(R)}\|
        \leq M\|\ket{V^{(N)}(R)}\|.\notag
    \end{align}
    The last expression tends to zero. The squeeze theorem proves the
    asserted complex limit.
\end{proof}

\begin{lemma}[Preservation of a unit overlap modulus]
    \label{thm:spc_aux_retained_overlap_one}
    \lean{MPSTensor.mpvOverlap_norm_tendsto_one_of_eventual_decomposition}
    \leanok
    \uses{def:mpv_state, def:mpv_overlap}
    Let $A,B,C,R$ have the same physical dimension and arbitrary finite
    bond dimensions. Suppose $\|\ket{V^{(N)}(A)}\|\leq M$ eventually
    for some $M\in\R$, and
    $\ket{V^{(N)}(B)}=\ket{V^{(N)}(C)}+\ket{V^{(N)}(R)}$ eventually.
    Assume $\|\ket{V^{(N)}(R)}\|\to0$.
    Convergence $|O_{AB}(N)|\to1$ then implies $|O_{AC}(N)|\to1$.
    Neither positivity nor a relation between the bond dimensions is
    required; the decomposition need not hold at length zero.
\end{lemma}

\begin{proof}\leanok
    \uses{thm:spc_aux_overlap_remainder_zero}
    The reverse triangle inequality gives
    \begin{align}
        \bigl||O_{AB}(N)|-|O_{AC}(N)|\bigr|
         & \leq|O_{AB}(N)-O_{AC}(N)|.\notag
    \end{align}
    The right-hand side tends to zero by
    Lemma~\ref{thm:spc_aux_overlap_remainder_zero}; hence the two overlap
    moduli have the same limit.
\end{proof}

\begin{lemma}[Irreducibility of a uniquely stationary support compression]
    \label{thm:spc_aux_supported_compression_irreducible}
    \lean{MPSTensor.isIrreducibleMap_of_unital_support_compression}
    \leanok
    Let $B$ be a tensor with bond dimension $k$, and suppose
    $\sigma\geq0$, $\tr\sigma=1$, and $\mathcal E_B^*(\sigma)=\sigma$.
    Suppose $\sigma$ is the unique trace-one matrix fixed by
    $\mathcal E_B^*$, with uniqueness among all such complex matrices,
    not merely positive semidefinite ones.
    Let $K:\C^D\to\C^k$ be an isometry with $D>0$ and
    $KK^\dagger=\operatorname{supp}\sigma$, and put
    $A^i=K^\dagger B^iK$. If $A$ is unital, then its primal transfer
    map $\mathcal E_A$ is irreducible.
    This is an auxiliary preparation result for
    \cite[Appendix~C, lines~2653--2717]{Schuch2011Phases}, using the
    channel criterion of \cite[Theorem~6.3]{Wolf2012Quantum}.
\end{lemma}

\begin{proof}\leanok
    \uses{thm:spt_stationary_density_compression}
    Compression to the entire support makes
    $\rho=K^\dagger\sigma K$ positive definite. Stationary support
    compression preserves its fixed-point equation, trace one, and
    uniqueness among all trace-one fixed matrices.
    Unitality of $A$ makes $\mathcal E_A^*$ a quantum channel.
    If $Z$ is any fixed matrix of that channel, then
    $Z+(1-\tr Z)\rho$ is a trace-one fixed matrix. Uniqueness gives
    $Z+(1-\tr Z)\rho=\rho$, hence $Z=(\tr Z)\rho$.
    A channel with a faithful fixed point and a one-dimensional fixed
    space is irreducible. Finally, irreducibility of a positive matrix
    map is equivalent to irreducibility of its trace adjoint, so
    $\mathcal E_A$ is irreducible as well.
\end{proof}

\subsection{Reduction to the stationary support}

\begin{lemma}[Adjoint transfer along a stationary support]
    \label{thm:spc_aux_adjoint_support_intertwining}
    \lean{MPSTensor.adjointMap_compression_lift}
    \leanok
    \uses{def:transfer_map}
    Let $B$ be a tensor and let $\sigma\geq0$ satisfy
    $\mathcal E_B^*(\sigma)=\sigma$. Let $K:\C^D\to\C^k$ be an isometry
    onto the support of $\sigma$, and put $A^i=K^\dagger B^iK$.
    Then, for every $D\times D$ matrix $Z$,
    \begin{align}
        \mathcal E_B^*(KZK^\dagger)
         & =K\mathcal E_A^*(Z)K^\dagger.\notag
    \end{align}
    Neither unitality nor normalization of $\sigma$ is required.
\end{lemma}

\begin{proof}\leanok
    \uses{def:transfer_map}
    Stationarity makes the support invariant under every adjoint letter.
    Thus the adjoint transfer of $KZK^\dagger$ remains supported there.
    Compression gives $\mathcal E_A^*(Z)$, and expansion along the
    isometry recovers the original supported matrix.
\end{proof}


\begin{lemma}[Preparation on the full stationary support]
    \label{thm:spc_aux_stationary_support_preparation}
    \lean{MPSTensor.prepare_stationary_support_compression}
    \leanok
    Let $B$ be a unital tensor with bond dimension $k$.
    Suppose $\sigma\geq0$, $\tr\sigma=1$, and
    $\mathcal E_B^*(\sigma)=\sigma$, and suppose $\sigma$ is the unique
    trace-one complex matrix fixed by $\mathcal E_B^*$.
    Let $K:\C^D\to\C^k$ be an isometry satisfying
    $KK^\dagger=\operatorname{supp}\sigma$, and put
    $A^i=K^\dagger B^iK$ and $\rho=K^\dagger\sigma K$.
    Then $A$ is unital, $\mathcal E_A$ is irreducible, and
    $\rho>0$, $\tr\rho=1$, and $\mathcal E_A^*(\rho)=\rho$.
    In particular, $D>0$ follows from these hypotheses.
    This is an auxiliary preparation result for
    \cite[Appendix~C, lines~2653--2717]{Schuch2011Phases}, using
    \cite[Lemma~6.4 and Theorem~6.3]{Wolf2012Quantum}.
\end{lemma}

\begin{proof}\leanok
    \uses{thm:spt_stationary_density_compression,
        thm:spc_aux_supported_compression_irreducible}
    Write $P=KK^\dagger$. Stationarity gives
    $(I-P)(B^i)^\dagger P=0$, hence $PB^i(I-P)=0$.
    Thus $K^\dagger B^iP=K^\dagger B^i$, and
    \begin{align}
        \sum_i A^i(A^i)^\dagger
         & =K^\dagger\Bigl(\sum_i B^iP(B^i)^\dagger\Bigr)K
        =K^\dagger\Bigl(\sum_i B^i(B^i)^\dagger\Bigr)K=I.\notag
    \end{align}
    Compression to the full support makes $\rho$ positive definite;
    stationary support compression gives its trace and fixed-point equation.
    The irreducibility criterion for a uniquely stationary support compression
    now applies. The trace-one condition excludes an empty support.
\end{proof}

\begin{lemma}[Normality from an ambient transfer bound]
    \label{thm:spc_aux_stationary_support_normality}
    \lean{MPSTensor.isNormalTensor_of_unital_stationary_support_compression}
    \leanok
    \uses{def:nt_cpsv}
    Let $B$ be unital, and let $\sigma\geq0$ be its unique trace-one complex
    adjoint stationary matrix. Let $K:\C^D\to\C^k$ be an isometry onto
    the entire support of $\sigma$, and put $A^i=K^\dagger B^iK$.
    Suppose $q<1$ and every $\lambda\in\sigma(\mathcal E_B)$ with
    $\lambda\neq1$ satisfies $|\lambda|\leq q$.
    Then $A$ is a normal tensor: its Kraus family is irreducible,
    $r(\mathcal E_A)=1$, and its unit-modulus spectrum is $\{1\}$.
    The transfer bound is an explicit hypothesis.
    This auxiliary consequence uses the normal-tensor definition of
    \cite[lines~231--235]{Cirac2016MPDO_arXiv} in the support setting of
    \cite[Appendix~C, lines~2653--2717]{Schuch2011Phases}.
\end{lemma}

\begin{proof}\leanok
    \uses{thm:spc_aux_stationary_support_preparation,
        thm:spc_aux_adjoint_support_intertwining}
    Preparation gives unitality and irreducibility. Positivity and
    $\mathcal E_A(I)=I$ give spectral radius one by
    \cite[Proposition~6.1]{Wolf2012Quantum}.
    Stationary support invariance gives
    \begin{align}
        \mathcal E_B^*(KZK^\dagger)
         & =K\mathcal E_A^*(Z)K^\dagger.\notag
    \end{align}
    The expansion $Z\mapsto KZK^\dagger$ is injective, since compression
    recovers $Z$. Therefore every eigenvalue of $\mathcal E_A^*$ is an
    eigenvalue of $\mathcal E_B^*$. The bilinear trace adjoint has the
    same characteristic polynomial as the original endomorphism, so
    $\sigma(\mathcal E_A)\subseteq\sigma(\mathcal E_B)$.
    The bound $q<1$ excludes every peripheral eigenvalue other than one,
    while the identity has eigenvalue one.
\end{proof}


% Finite-dimensional auxiliaries for the stationary-support reduction.
% These entries do not identify a physical Hamiltonian gap with transfer contraction.

\begin{lemma}[Equality of transfer and adjoint fixed-space dimensions]
    \label{thm:spc_aux_adjoint_fixed_dimension}
    \lean{MPSTensor.finrank_adjointFixedSpace_eq_transferFixedSpace}
    \leanok
    \uses{def:transfer_map}
    For any finite-dimensional tensor $B$, the complex dimensions of the
    fixed spaces of $\E_B$ and its trace adjoint $\E_B^*$ agree:
    $\dim_{\C}\ker(\Id-\E_B^*)=\dim_{\C}\ker(\E_B-\Id)$.
    Positivity, normalization, irreducibility and positive bond dimension
    are not required. This is a finite-dimensional auxiliary identity for
    \cite[Appendix~C, lines~2653--2717]{Schuch2011Phases}.
\end{lemma}

\begin{proof}\leanok
    In the matrix-unit basis and the corresponding transposed basis,
    the matrix of the trace adjoint of $\Id-\E_B$ is the transpose of
    the matrix of $\Id-\E_B$. Transposition preserves rank.
    The rank--nullity identity therefore gives equal kernel dimensions;
    replacing $\Id-\E_B$ by its negative does not change its kernel.
\end{proof}

\begin{lemma}[Positive-length decomposition along a triangular support]
    \label{thm:spc_aux_supported_periodic_split}
    \lean{MPSTensor.mpv_eq_compression_add_transient_of_upperZero}
    \leanok
    \uses{def:mpv_state}
    Let $K:\C^D\to\C^k$ be an isometry, put $P=KK^\dagger$ and
    $Q=\Id-P$, and let $C^i=K^\dagger B^iK$.
    Suppose $PB^iQ=0$ for every physical letter $i$.
    With $R^i=QB^iQ$, every $N>0$ satisfies
    $\ket{V^{(N)}(B)}=\ket{V^{(N)}(C)}+\ket{V^{(N)}(R)}$.
    No normalization or nonzero bond dimension is required.
    This is the triangular trace decomposition used in
    \cite[proof of Theorem~TIcanonical, lines~785--815]{PerezGarcia2007Matrix},
    and in \cite[Appendix~C, lines~2653--2717]{Schuch2011Phases}.
\end{lemma}

\begin{proof}\leanok
    \uses{thm:spt_isometric_bond_restriction}
    Put $L^i=PB^iP$. Then $L^iK=KC^i$ and $PL^i=L^i$,
    so the supported isometric restriction gives
    $\tr(L^w)=\tr(C^w)$ for every nonempty word $w$.
    The triangular condition shows that removing the off-diagonal block
    does not change word traces. Since $L^iR^j=R^iL^j=0$, induction on
    the nonempty word gives $(L+R)^w=L^w+R^w$.
    Taking the trace proves the asserted decomposition.
    Length zero is excluded because its coefficient is the bond dimension.
\end{proof}

\begin{lemma}[Identification after removal of a vanishing periodic remainder]
    \label{thm:spc_aux_retained_dimension_gauge}
    \lean{MPSTensor.exists_dim_eq_gaugePhase_of_vanishing_periodic_remainder}
    \leanok
    \uses{def:irreducible_tensor, def:gauge_phase_equiv, def:mpv_overlap}
    Let $A,B,C,R$ have the same physical dimension, with arbitrary finite
    bond dimensions for $B,R$ and positive bond dimensions $D,E$ for $A,C$.
    Suppose $A,C$ are irreducible and admit positive definite adjoint
    stationary matrices $\sigma,\tau$.
    Suppose $\|\ket{V^{(N)}(A)}\|\leq M$ eventually for some $M\in\R$,
    and $\ket{V^{(N)}(B)}=\ket{V^{(N)}(C)}+\ket{V^{(N)}(R)}$ eventually.
    If $\|\ket{V^{(N)}(R)}\|\to0$ and $|O_{AB}(N)|\to1$, then $D=E$.
    After this identification there exist $X\in\GL_D(\C)$ and
    $\zeta\in\C\setminus\{0\}$ such that
    $C^i=\zeta XA^iX^{-1}$ for every $i$.
    This is a conditional auxiliary reduction for
    \cite[Appendix~C, lines~2653--2717]{Schuch2011Phases}, using
    \cite[Lemma~equalMPS]{Cirac2016MPDO_arXiv}.
\end{lemma}

\begin{proof}\leanok
    \uses{thm:spc_aux_retained_overlap_one, thm:tp_gauge_from_adjoint_fixed,
        thm:rect_overlap_norm_one_dim_eq, thm:nt_overlap_norm_one_gauge}
    Removing the vanishing remainder gives $|O_{AC}(N)|\to1$.
    The pure bond gauges defined by $\sigma^{1/2}$ and $\tau^{1/2}$
    make $A,C$ trace-preserving, preserve irreducibility, and leave
    all periodic coefficients unchanged.
    The rectangular overlap theorem first gives $D=E$.
    Its equal-dimension counterpart then gives gauge equivalence up to a
    nonzero scalar. Composing with the inverse normalization gauges
    proves the statement for the original tensors.
\end{proof}




\begin{lemma}[Identification of a stationary support by its limiting overlap]
    \label{thm:spc_aux_stationary_support_identification}
    \lean{MPSTensor.exists_dim_eq_gaugePhase_of_unital_stationary_support_overlap}
    \leanok
    \uses{def:nt_cpsv, def:gauge_phase_equiv, def:mpv_overlap}
    Let $A,B$ have positive bond dimension $D$ and finite bond dimension $k$,
    respectively, with the same physical dimension. Suppose $A$ is irreducible
    and has a positive definite adjoint stationary matrix $\tau$.
    Suppose $B$ is unital and $\dim_{\C}\ker(\E_B-\Id)=1$.
    Let $\sigma\geq0$ satisfy $\tr\sigma=1$ and $\E_B^*(\sigma)=\sigma$.
    Let $K:\C^E\to\C^k$ be an isometry with
    $KK^\dagger=\operatorname{supp}\sigma$, and put $C^i=K^\dagger B^iK$.
    Suppose $0\leq q<1$ and every $z\in\operatorname{spec}(\E_B)$ with
    $z\ne1$ satisfies $|z|\leq q$.
    If $\|\ket{V^{(N)}(A)}\|\leq M$ eventually for some $M\in\R$ and
    $|O_{AB}(N)|\to1$, then $C$ is normal and $D=E$.
    After identifying the bond dimensions, there exist $X\in\GL_D(\C)$ and
    $\zeta\in\C\setminus\{0\}$ with $C^i=\zeta XA^iX^{-1}$ for every $i$.
    The spectral bound and overlap limit are explicit assumptions in this
    auxiliary reduction for
    \cite[Appendix~C, lines~2653--2717]{Schuch2011Phases}.
\end{lemma}

\begin{proof}\leanok
    \uses{thm:spc_aux_adjoint_fixed_dimension,
        thm:spc_aux_stationary_support_preparation,
        thm:spc_aux_stationary_support_normality,
        thm:spc_aux_supported_periodic_split,
        thm:spt_compact_transient_mpv_decay,
        thm:spc_aux_retained_dimension_gauge}
    Trace duality makes the adjoint fixed space one-dimensional.
    Since $\tr\sigma=1$, that space is $\C\sigma$.
    Every trace-one adjoint fixed matrix is therefore $\sigma$.
    Preparation on the full stationary support yields an irreducible unital
    tensor $C$ with faithful adjoint stationary matrix $K^\dagger\sigma K$;
    the strict spectral bound makes $C$ normal.
    Put $P=KK^\dagger$ and $R^i=(\Id-P)B^i(\Id-P)$.
    Stationarity gives $PB^i(\Id-P)=0$, hence
    $\ket{V^{(N)}(B)}=\ket{V^{(N)}(C)}+\ket{V^{(N)}(R)}$ for every $N>0$.
    The complementary-corner spectral estimate gives
    $\|\ket{V^{(N)}(R)}\|\to0$.
    Identification after removal of this vanishing remainder gives
    $D=E$ and the asserted bond gauge and scalar.
\end{proof}


\subsection{Propagation on finite periodic chains}

% Auxiliary finite-volume results. These do not assert exponential clustering.
% Source context: Hastings--Koma, arXiv:math-ph/0507008, Appendix A, Theorem A.2.
% The unrestricted interaction theorem is not identified with these specializations.

\paragraph{Finite periodic chains.}
Let $\Lambda=\mathbb Z/N\mathbb Z$, with $N\geq1$, and let
$\mathcal H=(\mathbb C^d)^{\otimes N}$. For $X\subseteq\Lambda$, write
$\mathcal A_X$ for the algebra of operators supported on $X$. The oriented
bond indexed by $j$ has support $e_j=\{j,j+1\}$. Distinct oriented bonds
are counted separately, even when their supports agree. All operator norms
below are Hilbert-space operator norms. For an operator $H$, put
$\tau_t(A)=e^{itH}Ae^{-itH}$. The restricted commutator norm $C_B(X,t)$, its
local recursion, and the weighted integral comparison are stated for an
arbitrary finite set of sites in Section~\ref{sec:qc_lieb_robinson}.

\begin{definition}[Ring support distance]
    \label{aux:ring-support-distance}
    \lean{MPSPreparation.ringSupportDistance}\leanok
    For nonempty $X,Y\subseteq\Lambda$, let $d(X,Y)$ be the least integer $r\geq0$
    for which there exist $x\in X$, $y\in Y$, and $m\in\mathbb Z$ such that
    $|m|\leq r$ and $y=x+m\pmod N$.
\end{definition}

\begin{lemma}[Displacement characterization]
    \label{aux:ring-support-distance-separation}
    \lean{MPSPreparation.isSeparatedBy_iff_lt_ringSupportDistance}\leanok
    \uses{aux:ring-support-distance}
    For nonempty supports and $r\in\mathbb N$, no displacement of absolute value
    at most $r$ connects $X$ to $Y$ if and only if $r<d(X,Y)$.
\end{lemma}
\begin{proof}\leanok
    Use the defining minimum and its minimizing displacement.
\end{proof}

\begin{lemma}[Finite-ring Lieb--Robinson estimate]
    \label{aux:finite-ring-lieb-robinson}
    \lean{MPSPreparation.norm_heisenberg_commutator_le_exp_abs_of_disjoint}\leanok
    \uses{aux:restricted-commutator-recursion,aux:restricted-commutator-continuity,
        aux:restricted-initial-commutator-bound,aux:restricted-initial-commutator-vanishing,
        aux:ring-support-distance,aux:restricted-commutator-evaluation,
        aux:restricted-commutator-negative-time,aux:weighted-volterra-comparison}
    Let $H=\sum_j h_j$, with Hermitian $h_j\in\mathcal A_{e_j}$ and
    $\|h_j\|\leq1$. Suppose $A\in\mathcal A_X$ and $B\in\mathcal A_Y$,
    where $X,Y$ are nonempty and disjoint. For every $a\geq0$ and $t\in\mathbb R$,
    \[
        \|[\tau_t(A),B]\|
        \leq2\|A\|\|B\|\,|X|e^{-a d(X,Y)}
        \bigl(e^{8e^a|t|}-1\bigr).
    \]
    The constants are independent of $N$ and $d$. This is a nearest-neighbor
    finite-periodic specialization of Hastings--Koma, Appendix A, Theorem A.2;
    it is not the unrestricted interaction theorem.
\end{lemma}
\begin{proof}\leanok
    Weight each nonempty support $Z$ by $|Z|e^{-a d(Z,Y)}$.
    If $e_j$ meets $Z$, then $d(Z,Y)\leq d(e_j,Y)+1$.
    At most $2|Z|$ oriented bonds meet $Z$, and each bond has at most two sites.
    Thus the weighted row sum in the commutator recursion is at most $8e^a$.
    A finite-dimensional integral comparison gives the exponential factor.
    Vanishing of the initial commutator gives the subtraction of one. Apply the
    same argument to $-H$ for negative times, and multiply by $\|A\|$.
\end{proof}

\begin{lemma}[Contractivity of cyclic local embedding]
    \label{aux:cyclic-local-embedding-norm}
    \lean{MPOTensor.norm_embedLocalOperator_le}\leanok
    Embedding an operator on $L$ consecutive sites into an $N$-site ring,
    where $L\leq N$, does not increase its Hilbert-space operator norm.
\end{lemma}
\begin{proof}\leanok
    The embedding is a star algebra homomorphism and is therefore contractive.
\end{proof}

\begin{lemma}[Support under a site permutation]
    \label{aux:permuted-operator-support}
    \lean{QuantumCircuit.supportedOperators_submatrix_siteEquiv}\leanok
    A permutation $p$ of the sites sends an operator supported on $X$ to an
    operator supported on $p(X)$.
\end{lemma}
\begin{proof}\leanok
    Check the assertion for elementary tensor-product operators, then extend by
    linearity to their span.
\end{proof}

\begin{lemma}[Support of a cyclic two-site embedding]
    \label{aux:cyclic-two-site-embedding-support}
    \lean{MPSPreparation.embedLocalOperator_twoSite_mem_supportedOperators}\leanok
    \uses{aux:permuted-operator-support}
    For $N\geq2$, the cyclic embedding of a two-site operator at $j$ belongs to
    $\mathcal A_{e_j}$, including the bond crossing from the last site to the first.
\end{lemma}
\begin{proof}\leanok
    The initial two-site window has the stated support. Translate it around the
    ring and use covariance of support under a permutation of sites.
\end{proof}

\begin{lemma}[Translation-invariant finite-ring propagation bound]
    \label{aux:interaction-hamiltonian-lieb-robinson}
    \lean{MPSTensor.norm_interactionHamiltonian_commutator_le_exp_abs_of_disjoint}\leanok
    \uses{aux:finite-ring-lieb-robinson,aux:cyclic-local-embedding-norm,
        aux:cyclic-two-site-embedding-support}
    Let $h$ be a Hermitian two-site interaction with $\|h\|\leq1$, and let
    $H_N=\sum_j h_{j,j+1}$ on the periodic chain of length $N\geq2$.
    For nonempty disjoint supports $X,Y$ and $A\in\mathcal A_X$, $B\in\mathcal A_Y$,
    the estimate
    \[
        \|[e^{itH_N}Ae^{-itH_N},B]\|
        \leq2\|A\|\|B\|\,|X|e^{-a d(X,Y)}
        \bigl(e^{8e^a|t|}-1\bigr)
    \]
    holds for every $a\geq0$ and $t\in\mathbb R$.
\end{lemma}
\begin{proof}\leanok
    Cyclic embedding preserves Hermiticity, has the indicated bond support,
    and is contractive. Apply the preceding finite-ring estimate.
\end{proof}

\begin{lemma}[Gaussian-filtered finite-ring commutators]
    \label{lem:spt_finite_ring_gaussian_commutator}
    \lean{MPSTensor.interactionHamiltonian_isHermitian,
        MPSTensor.norm_gaussianFilteredTerm_interactionHamiltonian_commutator_le,
        MPSTensor.norm_gaussianFilteredTerm_interactionHamiltonian_commutator_le_exp}
    \leanok
    \discussion{9879}
    \uses{aux:interaction-hamiltonian-lieb-robinson, thm:hastings_filtered_locality}
    Let $H_N$ be the periodic Hamiltonian of a Hermitian two-site interaction
    of norm at most one. Let $A\in\mathcal A_X$ and $B\in\mathcal A_Y$ for
    nonempty disjoint $X,Y$. For $q,T,a\geq0$ and $\Delta\in\mathbb R$, let
    $\widetilde A$ be the Gaussian average of $e^{itH_N}Ae^{-itH_N}$ with
    variance $\max(q/\Delta^2,0)$. Then
    \begin{align}
        \|[\widetilde A,B]\|
         & \leq 2\|A\|\|B\|\,|X|e^{-a d(X,Y)}
        \bigl(e^{8e^aT}-1\bigr)
        +2\|A\|\|B\|e^{-(T\Delta)^2/(2q)}.
        \label{eq:spt_gaussian_commutator}
    \end{align}
    The constants are independent of the ring length and physical dimension.
    At $q=0$, the final exponential is read as $1$.
    This is the periodic specialization of the Gaussian locality estimate
    of \cite[equation \texttt{lbnd}]{Hastings2006Solving}.
    For $0<\ell\leq d(X,Y)$ and $\Delta>0$, the choice
    $T=\ell/(16e)$ and $q=T\Delta$ gives
    \begin{align}
        \|[\widetilde A,B]\|
         & \leq 2\|A\|\|B\|\,|X|e^{-\ell/2}
        +2\|A\|\|B\|e^{-\ell\Delta/(32e)}.
        \label{eq:spt_gaussian_commutator_exp}
    \end{align}
\end{lemma}
\begin{proof}\leanok
    On $|t|\leq T$, monotonicity replaces the time factor in the periodic
    propagation estimate by $e^{8e^aT}-1$. On the remaining times use
    $\|[e^{itH_N}Ae^{-itH_N},B]\|\leq2\|A\|\|B\|$.
    The Gaussian mass outside $[-T,T]$ is at most
    $e^{-(T\Delta)^2/(2q)}$. Splitting the integral gives
    (\ref{eq:spt_gaussian_commutator}).
    For the stated parameter choice, $8eT=\ell/2$ and
    $(T\Delta)^2/(2q)=\ell\Delta/(32e)$. Also
    $e^{-d(X,Y)}(e^{\ell/2}-1)\leq e^{-\ell/2}$, which yields
    (\ref{eq:spt_gaussian_commutator_exp}).
\end{proof}

\begin{lemma}[Uniform Gaussian ground estimates along an exact MPS path]
    \label{lem:spt_exact_ground_gaussian_bounds}
    \lean{MPSTensor.ExactMPSGroundPath.exists_uniform_ground_line_spectral_gap,
        MPSTensor.ExactMPSGroundPath.exists_uniform_gaussian_ground_bounds}
    \leanok
    \discussion{9879}
    \uses{def:spt_exact_mps_gapped_phase, thm:hastings_gaussian_filtered_spectral,
        thm:hastings_filtered_no_ground_component}
    Let a symmetric gapped interaction path have a continuous exact MPS
    realization of its unique ground lines. There is a single $\Delta>0$
    such that the following holds at every parameter and ring length
    $N\geq2$. Write $H_N$ for the Hamiltonian, $\psi_N$ for the nonzero
    periodic MPS vector and $P_N$ for the projection onto $\mathbb C\psi_N$.
    For every operator $W$ and $q\geq0$, its Gaussian average with variance
    $q/\Delta^2$ satisfies
    \begin{align}
        \|(1-P_N)\widetilde W\psi_N\|
         & \leq \|W\|e^{-q/2}\|\psi_N\|.
        \label{eq:spt_gaussian_excited_ground}
    \end{align}
    If $\langle\psi_N,W\psi_N\rangle=0$, then also
    \begin{align}
        \|\widetilde W\psi_N\|
         & \leq \Delta^{-1}e^{-q/2}\|[H_N,W]\psi_N\|.
        \label{eq:spt_gaussian_centered_ground}
    \end{align}
    These are the filtered ground estimates of
    \cite[equation \texttt{gs}]{Hastings2006Solving} and
    \cite[equations \texttt{tighter}, \texttt{tighter2}]{Hastings2007AreaLaw}
    with a constant supplied by the physical path.
\end{lemma}
\begin{proof}\leanok
    The least energies in the gap condition and the unique-ground-line
    condition agree: each is at most the other. Thus every eigenvalue is
    either the ground energy or at least $\Delta$ above it, and the ground
    eigenspace is $\mathbb C\psi_N$. Apply the Gaussian spectral estimate
    to obtain (\ref{eq:spt_gaussian_excited_ground}). For centered $W$,
    $W\psi_N$ is orthogonal to that whole line; the commutator estimate gives
    (\ref{eq:spt_gaussian_centered_ground}). No normalization of $\psi_N$
    is needed.
\end{proof}

\begin{definition}[The energy-lowering part of an observable]
    \label{def:spt_negative_frequency_part}
    \lean{TNLean.Lattice.negativeFrequencyPart}
    \leanok
    \discussion{9879}
    Let $H$ be a finite Hermitian Hamiltonian with an orthonormal
    eigenbasis $e_i$ and energies $E_i$. For an observable $A$, define
    \[
        A^- = \sum_{E_i<E_j}
        \langle e_i,Ae_j\rangle\,|e_i\rangle\langle e_j|.
    \]
    This is the sharp negative-frequency part of the Heisenberg
    evolution $e^{itH}Ae^{-itH}$ in
    \cite[Section 3, equations (3.1)--(3.6)]{Hastings2006Spectral}.
\end{definition}

\begin{lemma}[Exact separation of the connected correlation]
    \label{lem:spt_negative_frequency_ground_correlation}
    \lean{TNLean.Lattice.negativeFrequencyPart_apply_ground_eq_zero,
        TNLean.Lattice.inner_ground_projected_eq_negativeFrequency_commutator,
        TNLean.Lattice.inner_ground_connected_eq_negativeFrequency_commutator}
    \leanok
    \discussion{9879}
    \uses{def:spt_negative_frequency_part}
    Suppose $E_0\le E_i$ for every $i$, let $\psi$ belong to the
    eigenspace at $E_0$, and let $P_0$ project onto that full eigenspace.
    Then $A^-\psi=0$ and
    \[
        \langle\psi,A(1-P_0)B\psi\rangle
        =\langle\psi,[A^-,B]\psi\rangle.
    \]
    If the ground eigenspace is $\mathbb C\psi$, this is
    \[
        \langle\psi,AB\psi\rangle-
        \frac{\langle\psi,A\psi\rangle\langle\psi,B\psi\rangle}
        {\|\psi\|^2}
        =\langle\psi,[A^-,B]\psi\rangle.
    \]
\end{lemma}
\begin{proof}\leanok
    An energy-lowering matrix element cannot act nontrivially on a
    ground vector. Expand an arbitrary vector orthogonal to the ground
    eigenspace in the energy basis. Between that vector and a ground
    bra, every nonzero matrix element of $A$ lowers energy, so it is
    retained by $A^-$. Apply this to $(1-P_0)B\psi$.
    The term $BA^-\psi$ vanishes. In a one-dimensional ground space,
    the rank-one projection formula gives the last identity without
    choosing a normalized vector.
\end{proof}

\begin{lemma}[Negative-frequency identity along the actual MPS path]
    \label{cor:spt_exact_ground_negative_frequency_identity}
    \lean{MPSTensor.ExactMPSGroundPath.negativeFrequency_connected_identity}
    \leanok
    \discussion{9879}
    \uses{lem:spt_exact_ground_gaussian_bounds,lem:spt_negative_frequency_ground_correlation}
    At each parameter and every ring length $N\ge2$ of an exact MPS
    ground path, the ordinary connected correlation of its actual
    periodic vector equals the actual negative-frequency commutator
    above. No transfer, correlation-decay or continuous-gauge input is
    required.
\end{lemma}
\begin{proof}\leanok
    The physical path supplies the least ground energy and the actual
    full ground line. Apply the preceding lemma to its Hamiltonian
    and periodic MPS vector.
\end{proof}

The sharp energy-lowering part need not be local. The following lemmas
replace it by a Gaussian regularization, control that regularization by a
time integral of local commutators, and derive ordinary exponential
clustering.

\begin{lemma}[The Gaussian negative-frequency selector]
    \label{lem:spt_gaussian_negative_frequency_cutoff}
    \lean{TNLean.Lattice.gaussianNegativeFrequencyCutoff,
        TNLean.Lattice.gaussianNegativeFrequencyCutoff_nonneg,
        TNLean.Lattice.gaussianNegativeFrequencyCutoff_le_one,
        TNLean.Lattice.gaussianNegativeFrequencyCutoff_le_exp_of_gap,
        TNLean.Lattice.abs_one_sub_gaussianNegativeFrequencyCutoff_le_exp_of_gap}
    \leanok
    \discussion{9879}
    For $q,\Delta>0$, let $X$ be a centered Gaussian of variance
    $\Delta^2/q$ and put $F_{q,\Delta}(\omega)=\mathbb P(X\ge\omega)$.
    Then $0\le F_{q,\Delta}(\omega)\le1$ and
    \[
        \omega\ge\Delta\ \Longrightarrow\
        F_{q,\Delta}(\omega)\le e^{-q/2},
        \qquad
        \omega\le-\Delta\ \Longrightarrow\
        |1-F_{q,\Delta}(\omega)|\le e^{-q/2}.
    \]
    This is the actual Gaussian selector of
    \cite[Lemma 3.1, equation (3.7)]{Hastings2006Spectral}, with
    $\alpha=\Delta^2/(2q)$.
\end{lemma}
\begin{proof}\leanok
    The moment-generating function is
    $\mathbb E(e^{tX})=\exp(\Delta^2t^2/(2q))$.
    Chernoff's bound at $\Delta$ is $e^{-q/2}$. The positive-frequency
    statement follows by inclusion of the corresponding tail events.
    The complementary event for $\omega\le-\Delta$ is contained in
    $\{-X\ge\Delta\}$, and $-X$ has the same sub-Gaussian bound.
\end{proof}

The following regularized time-integral argument identifies this actual
Gaussian selector. The operator locality estimates are proved after it.

\begin{definition}[The regularized Gaussian--Cauchy kernel]
    \label{def:spt_gaussian_cauchy_kernel}
    \lean{TNLean.Lattice.gaussianCauchyKernel}
    \leanok
    \discussion{9879}
    For $\alpha,\varepsilon>0$, set
    \[
        K_{\alpha,\varepsilon}(t)=
        \frac{e^{-\alpha t^2}}{2\pi(\varepsilon-it)}.
    \]
    This is the regularized negative-frequency time kernel of
    \cite[Lemma 3.1, equation (3.8)]{Hastings2006Spectral}.
\end{definition}

\begin{lemma}[Laplace representation and integrability of the actual kernel]
    \label{lem:spt_gaussian_cauchy_integrability}
    \lean{TNLean.Lattice.cauchyDenominator_ne_zero,
        TNLean.Lattice.norm_inv_cauchyDenominator_le,
        TNLean.Lattice.inv_cauchyDenominator_eq_integral,
        TNLean.Lattice.norm_gaussianCauchyKernel_le,
        TNLean.Lattice.integrable_gaussianCauchyKernel_mul_exp,
        TNLean.Lattice.norm_integral_gaussianCauchyKernel_mul_exp_le}
    \leanok
    \discussion{9879}
    \uses{def:spt_gaussian_cauchy_kernel}
    For $\alpha,\varepsilon>0$ and $\omega,t\in\mathbb R$,
    \[
        \frac1{\varepsilon-it}
        =\int_0^\infty e^{(-\varepsilon+it)s}\,ds,
        \qquad
        |K_{\alpha,\varepsilon}(t)|
        \le\frac{e^{-\alpha t^2}}{2\pi\varepsilon}.
    \]
    The function $K_{\alpha,\varepsilon}(t)e^{i\omega t}$ is absolutely
    integrable and
    \[
        \left|\int_{\mathbb R}
        K_{\alpha,\varepsilon}(t)e^{i\omega t}\,dt\right|
        \le \frac1{2\pi\varepsilon}\sqrt{\frac\pi\alpha}.
    \]
\end{lemma}
\begin{proof}\leanok
    The complex exponential integral converges since its exponent has
    real part $-\varepsilon<0$. Its antiderivative gives the Laplace
    identity. The denominator has norm at least $\varepsilon$, so
    the kernel is dominated by the displayed Gaussian. The oscillatory
    exponential has norm one; Gaussian integrability and the Gaussian
    integral formula finish the proof.
\end{proof}

The bound retains its dependence on $\varepsilon$. It is not a uniform
regularizer-limit estimate. The next lemma interchanges the two integrals,
identifies the Gaussian selector, and passes to the regularizer limit.

\begin{lemma}[The actual scalar Gaussian--Cauchy limit]
    \label{lem:spt_gaussian_cauchy_limit}
    \lean{TNLean.Lattice.gaussianCauchyTransform,
        TNLean.Lattice.gaussianCauchyTransform_eq_integral,
        TNLean.Lattice.tendsto_gaussianCauchyTransform_zero_regularizer,
        GaussianFilter.halfLine_gaussian_integral_eq_measureReal,
        TNLean.Lattice.tendsto_normalized_gaussianCauchyTransform,
        TNLean.Lattice.tendsto_normalized_gaussianCauchyTransform_cutoff}
    \leanok
    \discussion{9879}
    \uses{lem:spt_gaussian_cauchy_integrability,lem:spt_gaussian_negative_frequency_cutoff}
    Let $\nu_h$ be the centered Gaussian probability measure of variance
    $h>0$, and let $c_h$ be the density at zero of $\nu_{1/h}$.
    For $\varepsilon>0$, define the actual transform
    \[
        C_{h,\varepsilon}(\omega)=\int_{\mathbb R}
        \frac{e^{i\omega t}}{\varepsilon-it}\,d\nu_h(t).
    \]
    Then
    \[
        C_{h,\varepsilon}(\omega)=
        \int_0^\infty e^{-\varepsilon s}
        e^{-h(\omega+s)^2/2}\,ds,
        \qquad
        c_hC_{h,\varepsilon}(\omega)
        \longrightarrow \nu_{1/h}([\omega,\infty))
    \]
    as $\varepsilon\downarrow0$.
    In particular, $h=q/\Delta^2$ gives the actual Gaussian selector
    $F_{q,\Delta}$ above and its explicit gap-error bounds.
\end{lemma}
\begin{proof}\leanok
    Insert the half-line Laplace representation of the Cauchy
    denominator. Under $\nu_h\otimes ds$ on $s>0$, the absolute value of
    the joint integrand is $e^{-\varepsilon s}$, which is integrable.
    Fubini and the exact Gaussian characteristic function give the first
    identity. For the limit, the translated Gaussian
    $e^{-h(\omega+s)^2/2}$ is an integrable majorant independent of
    $\varepsilon>0$. Apply dominated convergence.
    The actual density formula for $\nu_{1/h}$, translation invariance
    of Lebesgue integration, and equality of open and closed integral
    endpoints give the normalization to the probability tail.
\end{proof}

The scalar frequency selector and its integral representation supply
the regularized observable, ground-vector comparison, and locality
argument below.

\begin{definition}[The actual regularized negative-frequency observable]
    \label{def:spt_regularized_negative_frequency_operator}
    \lean{TNLean.Lattice.gaussianNegativeFrequencyIntegrand,
        TNLean.Lattice.gaussianNegativeFrequencyOperator}
    \leanok
    \discussion{9879}
    \uses{lem:spt_gaussian_cauchy_limit}
    For a Hermitian $H$, an observable $A$, Gaussian time variance $h$,
    and positive regularizer $\varepsilon$, put
    \[
        A^-_{h,\varepsilon}
        =c_h\int_{\mathbb R}
        \frac{e^{itH}Ae^{-itH}}{\varepsilon-it}\,d\nu_h(t).
    \]
\end{definition}

\begin{lemma}[Actual regularized energy matrix elements]
    \label{lem:spt_regularized_frequency_matrix_elements}
    \lean{TNLean.Lattice.continuous_gaussianNegativeFrequencyIntegrand,
        TNLean.Lattice.norm_gaussianNegativeFrequencyIntegrand_le,
        TNLean.Lattice.integrable_gaussianNegativeFrequencyIntegrand,
        TNLean.Lattice.norm_gaussianNegativeFrequencyOperator_le,
        TNLean.Lattice.dotProduct_integral_gaussianNegativeFrequencyIntegrand_mulVec,
        TNLean.Lattice.dotProduct_gaussianNegativeFrequencyOperator_mulVec,
        TNLean.Lattice.inner_gaussianNegativeFrequencyOperator,
        TNLean.Lattice.tendsto_dotProduct_gaussianNegativeFrequencyOperator}
    \leanok
    \discussion{9879}
    \uses{def:spt_regularized_negative_frequency_operator,lem:spt_gaussian_cauchy_limit}
    The actual regularized observable is integrable and
    $\|A^-_{h,\varepsilon}\|\le c_h\varepsilon^{-1}\|A\|$.
    For energy eigenvectors $Hv=Ev$ and $Hw=E'w$,
    \[
        \langle v,A^-_{h,\varepsilon}w\rangle
        =c_hC_{h,\varepsilon}(E-E')\langle v,Aw\rangle.
    \]
    If $h>0$, the actual entries therefore converge as
    $\varepsilon\downarrow0$ to
    \[
        \nu_{1/h}([E-E',\infty))\langle v,Aw\rangle.
    \]
\end{lemma}
\begin{proof}\leanok
    Unitary conjugation preserves the norm, and the rational factor has
    norm at most $\varepsilon^{-1}$, so the finite Gaussian measure
    gives integrability. Between energy eigenvectors the actual
    Heisenberg evolution contributes $e^{it(E-E')}$. Passing the matrix
    element through the integral gives the actual scalar transform.
    Apply its proved Gaussian probability-tail limit.
\end{proof}

\begin{lemma}[Uniform locality after cancelling the Cauchy singularity]
    \label{lem:spt_regularized_frequency_uniform_locality}
    \lean{TNLean.Lattice.norm_commutator_gaussianNegativeFrequencyOperator_le}
    \leanok
    \discussion{9879}
    \uses{def:spt_regularized_negative_frequency_operator}
    Suppose $C,v\ge0$ and the actual dynamics satisfy
    $\|[A(t),B]\|\le C(e^{v|t|}-1)$ for all real $t$.
    For $T,\varepsilon>0$,
    \[
        \|[A^-_{h,\varepsilon},B]\|
        \le c_h\left(Cve^{vT}
        +\frac{2\|A\|\|B\|}{T}e^{-T^2/(2h)}\right).
    \]
    In particular the bound is independent of the positive regularizer.
\end{lemma}
\begin{proof}\leanok
    The rational factor is at most $|t|^{-1}$ for $t\ne0$.
    For $|t|\le T$, the inequality $e^x-1\le xe^x$ cancels this
    factor against the actual short-time increment. At $t=0$ the
    commutator vanishes. Outside the time window, the rational factor
    is at most $T^{-1}$ and the commutator norm at most
    $2\|A\|\|B\|$. Split the actual Gaussian integral and apply its
    Gaussian tail bound. No bound growing as $\varepsilon^{-1}$ is
    used for this passage.
\end{proof}

\begin{lemma}[Actual finite-ring regularized-frequency locality]
    \label{cor:spt_finite_ring_regularized_frequency_locality}
    \lean{MPSTensor.norm_negativeFrequency_interactionHamiltonian_commutator_le}
    \leanok
    \discussion{9879}
    \uses{lem:spt_regularized_frequency_uniform_locality,aux:interaction-hamiltonian-lieb-robinson}
    For an actual periodic nearest-neighbour Hamiltonian with Hermitian
    interaction of norm at most one and disjoint nonempty supports $X,Y$,
    the preceding bound holds with
    \[
        C=2\|A\|\|B\|\,|X|e^{-a d(X,Y)},
        \qquad v=8e^a,\qquad a\ge0.
    \]
    There is no supplied propagation or correlation estimate.
\end{lemma}
\begin{proof}\leanok
    The actual finite-ring Lieb--Robinson bound supplies exactly the
    required short-time increment. Apply the preceding lemma.
\end{proof}

\begin{lemma}[Actual Gaussian ground-vector separation]
    \label{lem:spt_gaussian_ground_frequency_separation}
    \lean{TNLean.Lattice.exists_gaussianNegativeFrequency_ground_limits,
        TNLean.Lattice.exists_gaussianNegativeFrequency_ground_commutator_limit}
    \leanok
    \discussion{9879}
    \uses{lem:spt_regularized_frequency_matrix_elements,lem:spt_gaussian_negative_frequency_cutoff}
    Suppose the full ground eigenspace of $H$ is $\mathbb C\psi$,
    every other energy is at least $\Delta>0$ above it, and
    $\langle\psi,A\psi\rangle=0$. For $q>0$, put $h=q/\Delta^2$.
    The vectors $A^-_{h,\varepsilon}\psi$ and
    $(A^-_{h,\varepsilon})^\dagger\psi$ have limits $\xi,\eta$ with
    \[
        \|\xi\|\le e^{-q/2}\|A\psi\|,
        \qquad \|A^\dagger\psi-\eta\|
        \le e^{-q/2}\|A^\dagger\psi\|.
    \]
    The scalar commutator therefore converges to a number $z$ satisfying
    \[
        |\langle\psi,AB\psi\rangle-z|
        \le 2e^{-q/2}\|A\|\|B\|\|\psi\|^2.
    \]
    No normalization of $\psi$ or centering of $B$ is required.
\end{lemma}
\begin{proof}\leanok
    Expand in an orthonormal energy basis and apply the matrix-element
    limit. Centering removes all ground-energy coefficients. The two
    probability-tail estimates bound the remaining coefficients, and
    Parseval gives the vector bounds. Move the adjoint across the inner
    product and apply Cauchy--Schwarz to the two error terms.
\end{proof}

\begin{lemma}[Centering an observable at an unnormalized vector]
    \label{lem:spt_actual_observable_centering}
    \lean{TNLean.Lattice.normalizedVectorExpectation,
        TNLean.Lattice.centeredObservable,
        TNLean.Lattice.norm_normalizedVectorExpectation_le,
        TNLean.Lattice.inner_centeredObservable_eq_zero,
        TNLean.Lattice.norm_centeredObservable_le,
        TNLean.Lattice.inner_centeredObservable_mul,
        TNLean.Lattice.centeredObservable_mem,
        TNLean.Lattice.centeredObservable_commutator}
    \leanok
    \discussion{9879}
    For $\psi\ne0$, write
    $\langle A\rangle_\psi=\langle\psi,A\psi\rangle/\|\psi\|^2$
    and $A_0=A-\langle A\rangle_\psi\Id$. Then
    $|\langle A\rangle_\psi|\le\|A\|$,
    $\langle\psi,A_0\psi\rangle=0$, and $\|A_0\|\le2\|A\|$.
    The centered observable gives the connected correlation:
    \[
        \frac{\langle\psi,A_0B\psi\rangle}{\|\psi\|^2}
        =\langle AB\rangle_\psi
        -\langle A\rangle_\psi\langle B\rangle_\psi.
    \]
    Subtracting this scalar identity preserves every linear space of
    observables containing $A$ and $\Id$, as well as the commutator with $B$.
\end{lemma}
\begin{proof}\leanok
    Use Cauchy--Schwarz, the operator-norm bound, and the triangle
    inequality. Expand the scalar identity and use that $\|\psi\|^2>0$.
\end{proof}

\begin{lemma}[Ordinary centered correlation from Gaussian frequency separation]
    \label{lem:spt_gaussian_ordinary_centered_correlation}
    \lean{TNLean.Lattice.norm_centered_ground_correlation_le_gaussian}
    \leanok
    \discussion{9879}
    \uses{lem:spt_gaussian_ground_frequency_separation,lem:spt_regularized_frequency_uniform_locality}
    Under the preceding ground-line and propagation hypotheses, for
    $q,T>0$ and $h=q/\Delta^2$,
    \[
        |\langle\psi,AB\psi\rangle|
        \le\left(2e^{-q/2}\|A\|\|B\|
        +c_h\left(Cve^{vT}
            +\frac{2\|A\|\|B\|}{T}e^{-T^2/(2h)}\right)\right)\|\psi\|^2.
    \]
\end{lemma}
\begin{proof}\leanok
    The regularized commutator is bounded uniformly in its positive
    regularizer. Pass this bound to its scalar limit and add the proved
    Gaussian ground-vector error.
\end{proof}

\begin{lemma}[Ordinary connected Gaussian bound for the physical path]
    \label{cor:spt_exact_ground_connected_gaussian_bound}
    \lean{MPSTensor.ExactMPSGroundPath.exists_uniform_centered_correlation_gaussian_bound,
        MPSTensor.ExactMPSGroundPath.exists_uniform_connected_correlation_gaussian_bound}
    \leanok
    \discussion{9879}
    \uses{lem:spt_gaussian_ordinary_centered_correlation,lem:spt_actual_observable_centering,
        lem:spt_exact_ground_gaussian_bounds,aux:interaction-hamiltonian-lieb-robinson}
    Along an exact MPS ground-state path there is one $\Delta>0$ such that,
    at every parameter and ring length $N\ge2$, arbitrary observables
    supported on disjoint nonempty sets $X,Y$ satisfy
    \[
        |\langle AB\rangle_\psi
        -\langle A\rangle_\psi\langle B\rangle_\psi|
        \le2\|A\|\|B\|
        \left(2e^{-q/2}+c_h\left(
            2|X|e^{-a d(X,Y)}ve^{vT}
            +\frac2T e^{-T^2/(2h)}\right)\right),
    \]
    where $q,T>0$, $a\ge0$, $h=q/\Delta^2$, and $v=8e^a$.
\end{lemma}
\begin{proof}\leanok
    The actual physical gap supplies the full ground line and the common
    spectral gap. The actual finite-ring propagation estimate supplies
    the time increment. Center the first observable at the actual
    nonzero periodic MPS vector, preserve its support, and divide the
    preceding estimate by the squared vector norm.
\end{proof}

\begin{lemma}[An exponential rate for the Gaussian frequency bound]
    \label{lem:spt_gaussian_frequency_clustering_rate}
    \lean{GaussianFilter.kernel_inv_zero_eq_sqrt,
        TNLean.Lattice.exists_gaussianFrequencyClusteringRate}
    \leanok
    \discussion{9879}
    If $\Delta,v>0$, there are $C,\kappa>0$ such that for $m\ge0$ and
    $\ell\ge1$,
    \[
        2e^{-\ell\Delta/(4v)}+
        \sqrt{\frac{\ell}{4\pi v\Delta}}
        \left(2mv e^{-\ell/2}
        +\frac{4v}{\ell}e^{-\ell\Delta/(4v)}\right)
        \le C(m+1)e^{-\kappa\ell}.
    \]
\end{lemma}
\begin{proof}\leanok
    The density normalization is $c_h=\sqrt{h/(2\pi)}$.
    Take $\kappa=\min(1/4,\Delta/(8v))$ and absorb the square-root
    factors using $\sqrt\ell\le\ell$ and
    $\ell e^{-\ell/2}\le4e^{-\ell/4}$.
\end{proof}

\begin{lemma}[Ordinary connected exponential clustering along the physical path]
    \label{thm:spt_exact_ground_connected_exponential_clustering}
    \lean{MPSTensor.ExactMPSGroundPath.exists_uniform_connected_exponential_clustering}
    \leanok
    \discussion{9879}
    \uses{cor:spt_exact_ground_connected_gaussian_bound,lem:spt_gaussian_frequency_clustering_rate}
    Along an exact MPS ground-state path there are $K,\kappa>0$,
    independent of the parameter and ring length $N\ge2$, such that
    observables supported on disjoint nonempty $X,Y$ obey
    \[
        |\langle AB\rangle_\psi
        -\langle A\rangle_\psi\langle B\rangle_\psi|
        \le K(|X|+1)\|A\|\|B\|e^{-\kappa\ell}
    \]
    whenever $1\le\ell\le d(X,Y)$.
\end{lemma}
\begin{proof}\leanok
    Set $a=1$, $v=8e$, $T=\ell/(2v)$, and $q=T\Delta$ in the
    preceding Gaussian bound. Both Gaussian errors are
    $e^{-\ell\Delta/(4v)}$, while the short-time term is bounded
    by $2|X|v e^{-\ell/2}$. Apply the exponential-rate lemma.
\end{proof}

\begin{lemma}[Normalized correlations depend only on the periodic rays]
    \label{lem:spt_positive_ray_expectations}
    \lean{TNLean.Lattice.normalizedVectorExpectation_smul,
        MPSTensor.mpvExpectation_eq_normalizedVectorExpectation,
        MPSTensor.SamePositiveMpvRay.mpvExpectation_eq,
        MPSTensor.SamePositiveMpvRay.mpvConnectedCorrelator_eq}
    \leanok
    \discussion{9879}
    \uses{lem:spt_actual_observable_centering}
    If two tensors generate the same periodic vector line at every
    positive length, all their normalized finite expectations and
    connected window correlators coincide. Their bond dimensions and
    scalar normalizations may differ.
\end{lemma}
\begin{proof}\leanok
    Equal singleton spans differ by an invertible scalar, including when
    both vectors vanish. The scalar and its complex conjugate cancel in
    the ratio defining the normalized expectation. Apply this identity
    to the three expectations defining the connected correlator.
\end{proof}

\begin{lemma}[Physical clustering for the actual two-window correlator]
    \label{cor:spt_physical_window_clustering}
    \lean{MPSTensor.ExactMPSGroundPath.exists_uniform_window_correlation_bound,
        MPSTensor.ExactMPSGroundPath.exists_uniform_ray_window_correlation_bound}
    \leanok
    \discussion{9879}
    \uses{thm:spt_exact_ground_connected_exponential_clustering,lem:spt_positive_ray_expectations}
    There are common $K,\kappa>0$ such that for any tensor representing
    a periodic ray family along the physical path, length-$L$ observables
    $X,Y$ with $L>0$, and $n\ge t\ge1$, their connected correlator on
    the ring of length $L+t+L+n$ satisfies
    \[
        |G_{L+t+L+n}(X,Y;L+t+1)|
        \le K(L+1)\|X\|\|Y\|e^{-\kappa t}.
    \]
\end{lemma}
\begin{proof}\leanok
    The two windows have at least $t$ sites between them in both cyclic
    directions, so their shortest-path distance is greater than $t$.
    Window embeddings are contractive and the first support has at most
    $L$ sites. Apply physical clustering and the preceding ray identity.
\end{proof}

\begin{lemma}[Actual finite-ring correlations converge to the limiting contraction]
    \label{lem:spt_actual_finite_ring_correlation_limit}
    \lean{MPSTensor.tendsto_mpvConnectedCorrelator_limitCorrelator,
        MPSTensor.norm_limitCorrelator_le_of_eventually_mpvConnectedCorrelator_le,
        MPSTensor.norm_transfer_spectrum_le_of_finite_mpvConnectedCorrelator_bound}
    \leanok
    \discussion{9879}
    \uses{def:physical_connected_correlator,lem:spt_transfer_rate_from_correlations}
    Let $A$ be trace preserving and injective after blocking $L>0$ sites,
    with a faithful stationary density $\rho$ of trace one. For fixed
    $t\ge1$ and length-$L$ observables $X,Y$, the actual finite connected
    correlator on $L+t+L+n$ sites converges, as $n\to\infty$, to the
    limiting fixed-point contraction. Every eventual bound on the
    finite correlator therefore bounds this contraction. In particular,
    a common eventual rate $r>0$ for all such observables bounds every
    transfer eigenvalue different from one by $r$.
\end{lemma}
\begin{proof}\leanok
    Block injectivity gives pointwise primitivity. The complementary
    transfer operator has spectral radius less than one, so its powers
    tend to zero. Substitute this in the exact finite-ring contraction
    formula; its denominator tends to one. Pass the finite bound through
    the norm limit and use the transfer eigenvalue observables.
    The auxiliary decay rate may depend on the tensor; no common transfer
    bound is used to prove convergence.
\end{proof}

\begin{lemma}[A common transfer bound for trace-preserving representatives]
    \label{thm:spt_physical_gap_tp_transfer_bound}
    \lean{MPSTensor.ExactMPSGroundPath.exists_uniform_tracePreserving_transfer_bound}
    \leanok
    \discussion{9879}
    \uses{cor:spt_physical_window_clustering,lem:spt_actual_finite_ring_correlation_limit}
    An exact physical MPS ground-state path has a common $0<r<1$ such
    that every trace-preserving block-injective representative of one of
    its periodic ray families, with a faithful stationary trace-one
    density, has all transfer eigenvalues different from one of modulus
    at most $r$.
\end{lemma}
\begin{proof}\leanok
    Take $r=e^{-\kappa}$ from the physical window bound. For each pair
    of observables the estimate holds for all return lengths $n\ge t$.
    Pass to the limiting contraction and then to the transfer spectrum.
\end{proof}

\begin{lemma}[An actual trace-preserving gauge with faithful fixed density]
    \label{lem:spt_injective_tp_gauge}
    \lean{MPSTensor.exists_posDef_fixedPoint_trace_one_of_isInjective_isTP,
        MPSTensor.exists_tpGauge_of_isInjective_unital}
    \leanok
    \discussion{9879}
    A trace-preserving one-site injective tensor of positive bond dimension
    has a faithful stationary density of trace one. Every unital one-site
    injective tensor $A$ of positive bond dimension has a bond gauge $B$
    that is trace preserving, one-site injective, has the same
    positive-length periodic rays and transfer spectrum, and has such a
    density.
\end{lemma}
\begin{proof}\leanok
    Primitive Perron--Frobenius theory supplies the positive-definite
    stationary density, which is normalized by its positive trace.
    For the gauge, derive the faithful adjoint fixed point $\Lambda$ of
    the unital transfer map and set
    $B^i=\Lambda^{1/2}A^i\Lambda^{-1/2}$.
    The adjoint fixed-point identity gives trace preservation. Invertible
    bond conjugation preserves injectivity, the periodic vectors, and
    the transfer spectrum.
\end{proof}

\begin{lemma}[A common unital transfer bound from the physical gap]
    \label{thm:spt_physical_gap_unital_transfer_bound}
    \lean{MPSTensor.ExactMPSGroundPath.exists_uniform_unital_transfer_bound}
    \leanok
    \discussion{9879}
    \uses{thm:spt_physical_gap_tp_transfer_bound,lem:spt_injective_tp_gauge}
    Along an exact physical MPS ground-state path there is a common
    $0<r<1$ bounding the modulus of every transfer eigenvalue different
    from one of every unital injective representative of a periodic ray
    family on the path. No transfer bound or canonical gauge is an
    additional hypothesis.
\end{lemma}
\begin{proof}\leanok
    Derive a trace-preserving gauge and its faithful fixed density.
    Apply the preceding physical transfer bound at blocking length one,
    and transport the spectrum back through the bond gauge.
\end{proof}

\begin{lemma}[Transfer spectral rate from limiting correlations]
    \label{lem:spt_transfer_rate_from_correlations}
    \lean{MPSTensor.norm_transfer_spectrum_le_of_limitCorrelator_bound}
    \leanok
    \discussion{9879}
    \uses{lem:ldp_limit_correlator_eigenvalue, def:physical_connected_correlator}
    Let the products of $L$ letters of $A$ span the bond matrix algebra,
    suppose $\sum_i(A^i)^\dagger A^i=\Id$, and let $\rho$ have trace one.
    Suppose $r>0$ and for every pair of observables $X,Y$ on $L$ sites
    there are $C\in\mathbb R$ and $n_0\in\mathbb N$ with
    $|C_{\mathrm{phys}}(X,Y;n)|\leq Cr^n$ for $n\geq n_0$.
    Every $\lambda\ne1$ in the transfer spectrum satisfies $|\lambda|\leq r$.
\end{lemma}
\begin{proof}\leanok
    The eigenvalue observables give $C_{\mathrm{phys}}(X,Y;n)=\lambda^n$. If
    $|\lambda|>r$, the ratio $|\lambda|/r$ is greater than one and its
    powers are unbounded. The assumed estimate at $n+n_0$ bounds these
    powers by $C$, a contradiction. The prefactors and starting
    separations need not be uniform over the observables.
\end{proof}

\subsection{Virtual classes at a compact limit}

\begin{theorem}[One virtual class from the actual physical gap]
    \label{thm:spt_exact_ground_physical_class}
    \lean{MPSTensor.ExactMPSGroundPath.exists_constant_virtual_class}
    \leanok
    \discussion{9879}
    \uses{thm:spt_physical_gap_unital_transfer_bound,
        thm:spc_aux_exact_ground_uniform_transfer_class}
    A symmetric gapped interaction path with a continuous exact MPS
    realization of its unique finite-chain ground lines has a common
    unit-modulus physical character and unital injective representatives
    of positive bond dimensions bounded by the ambient dimension.
    After removing the character, these representatives have unitary
    projective virtual actions whose factor systems all determine the
    same cohomology class. The minimal dimensions may change along the
    path, and no continuous choice of minimal tensors, stationary
    densities, or virtual matrices is required.
\end{theorem}
\begin{proof}\leanok
    Derive the common unital transfer bound from the physical gap and
    apply the compact-limit class-invariance theorem below.
\end{proof}

\begin{lemma}[Virtual class of an injective periodic ray]
    \label{lem:spt_periodic_ray_virtual_class}
    \lean{MPSTensor.bondDimension_eq_of_samePositiveMpvRay,
        MPSTensor.cohomologousTo_of_samePositiveMpvRay,
        MPSTensor.cohomologousTo_of_samePositiveMpvRay_character_covariance}
    \leanok
    \uses{thm:spt_positive_ray_finite_traces}
    Let $A$ and $B$ be injective tensors of positive bond dimensions with
    the same positive-length periodic rays. Their bond dimensions agree.
    Prescribed virtual actions for the same physical representation have
    equal cohomology classes, including when the actions remove separately
    chosen unit-modulus physical characters.
\end{lemma}
\begin{proof}\leanok
    The nonzero two- and three-site trace normalizations determine the
    common bond dimension and a scalar bond gauge. Transport the first
    virtual action through this gauge. On the second tensor, injective
    gauge-phase uniqueness forces the ratio of the physical characters
    to be one. Gauge independence then identifies the factor systems.
\end{proof}

\begin{lemma}[Prescribed classes at the endpoints of a physical path]
    \label{cor:spt_exact_ground_endpoint_class}
    \lean{MPSTensor.ExactMPSGroundPath.cohomologousTo_endpoints}
    \leanok
    \uses{thm:spt_exact_ground_physical_class, lem:spt_periodic_ray_virtual_class}
    Let $Q$ realize the unique finite-chain ground lines of a symmetric
    gapped interaction path. Suppose injective positive-dimensional
    tensors $A_0,A_1$ represent its endpoint periodic rays. Any prescribed
    virtual actions of $A_0,A_1$ for the common physical representation,
    after independent unit-modulus character changes, have equal
    cohomology classes.
\end{lemma}
\begin{proof}\leanok
    Compare each endpoint action with the corresponding action in the
    common class supplied by Theorem~\ref{thm:spt_exact_ground_physical_class}.
    Lemma~\ref{lem:spt_periodic_ray_virtual_class} identifies each pair.
\end{proof}

\begin{lemma}[Cohomology obstruction to an exact MPS gapped phase]
    \label{cor:spt_exact_phase_cohomology_obstruction}
    \lean{MPSTensor.IsSameExactMPSGappedPhase.cohomologousTo}
    \leanok
    \uses{def:spt_exact_mps_gapped_phase, cor:spt_exact_ground_endpoint_class}
    Let $A_0,A_1$ be one-site injective tensors of positive bond dimensions
    with prescribed virtual projective actions implementing their exact
    on-site symmetries. If they satisfy the independent exact MPS gapped
    phase condition, their virtual factor systems have equal cohomology
    classes. The common blocking, endpoint physical character changes,
    and orthogonal physical inclusions are those permitted by the phase
    condition.
\end{lemma}
\begin{proof}\leanok
    Blocking preserves each virtual factor system and injectivity.
    Each physical isometry also preserves injectivity. The physical
    intertwining relation transports the blocked virtual action to its
    included tensor after removing the endpoint character. Apply
    Lemma~\ref{cor:spt_exact_ground_endpoint_class} to the resulting
    endpoint representatives of the actual path.
\end{proof}

\begin{theorem}[Classification for injective endpoints with exact virtual covariance]
    \label{thm:spt_exact_phase_injective_classification}
    \lean{MPSTensor.isSameExactMPSGappedPhase_iff_cohomologous}
    \leanok
    \uses{cor:spt_exact_phase_cohomology_obstruction, thm:spt_exact_phase_forward}
    Let $A_0,A_1$ be one-site injective tensors of positive bond dimensions,
    with prescribed virtual projective representations implementing their
    exact on-site symmetries. They satisfy the independent exact MPS gapped
    phase condition if and only if the prescribed factor systems determine
    the same cohomology class. This is the injective, exact-symmetry scope
    of the construction in \cite[Section~II.F.2]{Schuch2011Phases}.
    Endpoints whose symmetry carries a scalar character are treated in
    Corollary~\ref{cor:spt_exact_phase_injective_character_classification}
    \cite{gap:spc11_spt_interpolation_upper_range}.
\end{theorem}
\begin{proof}\leanok
    The forward construction for cohomologous virtual actions gives the
    physical phase condition. The converse is
    Lemma~\ref{cor:spt_exact_phase_cohomology_obstruction}.
\end{proof}

\begin{lemma}[Endpoint character changes preserve the exact MPS gapped phase]
    \label{lem:spt_exact_phase_endpoint_character_change}
    \lean{MPSTensor.IsSameExactMPSGappedPhase.of_smul_left}
    \leanok
    \uses{def:spt_exact_mps_gapped_phase}
    Let $\chi$ be a unit-modulus character of $G$. If $A_0,A_1$ satisfy the
    independent exact MPS gapped phase condition for the physical actions
    $U^0,U^1$, they satisfy it for $\chi(g)U^0_g$ and $U^1_g$.
\end{lemma}
\begin{proof}\leanok
    The phase condition allows a unit-modulus character at each endpoint
    inclusion. After blocking $k$ sites, the $k$-fold tensor power of
    $\chi(g)U^0_g$ is $\chi(g)^kU^{0\otimes k}_g$. Multiplying the first
    endpoint character by $\chi^{-k}$ therefore keeps the same blocking,
    inclusions, interaction path and exact MPS realization.
\end{proof}

\begin{corollary}[Classification for injective endpoints with a symmetry character]
    \label{cor:spt_exact_phase_injective_character_classification}
    \lean{MPSTensor.isSameExactMPSGappedPhase_iff_cohomologous_of_character}
    \leanok
    \uses{thm:spt_exact_phase_injective_classification,
        lem:spt_exact_phase_endpoint_character_change}
    Let $A_0,A_1$ be one-site injective tensors of positive bond dimensions,
    let $\chi_0,\chi_1$ be unit-modulus characters of $G$, and let
    $V^0,V^1$ be virtual projective representations with
    \begin{align}
        \sum_j\bigl(U^p_g\bigr)_{ij}A_p^j
         & =\chi_p(g)\,V^p_{g^{-1}}A_p^i\bigl(V^p_{g^{-1}}\bigr)^{-1},
        \qquad p=0,1,
        \notag
    \end{align}
    for all $g$ and $i$, the tensor form of the symmetry relation with a
    character in \cite[Appendix~B]{Schuch2011Phases}. They satisfy the independent
    exact MPS gapped phase condition for $U^0,U^1$ if and only if the factor
    systems of $V^0$ and $V^1$ determine the same cohomology class.
\end{corollary}
\begin{proof}\leanok
    Following \cite[Section~II.C.2 and Appendix~B]{Schuch2011Phases},
    replace $U^p_g$ by the phase gauge $\chi_p(g)^{-1}U^p_g$. The virtual
    covariance then holds without a character, and
    Lemma~\ref{lem:spt_exact_phase_endpoint_character_change}, applied at
    both endpoints, shows that the phase condition is the same for the two
    gauges. Theorem~\ref{thm:spt_exact_phase_injective_classification}
    concludes.
\end{proof}

\begin{lemma}[The virtual class of an exact ground-state path under a transfer bound]
    \label{thm:spc_aux_exact_ground_uniform_transfer_class}
    \lean{MPSTensor.ExactMPSGroundPath.exists_constant_virtual_class_of_uniform_transfer_bound}
    \leanok
    Let $P$ be a symmetric gapped interaction path on the unit interval,
    with an independent continuous exact MPS realization $Q$ of its unique
    finite-chain ground lines. Suppose $0\le q<1$ has the following property:
    at every parameter, every positive-dimensional unital injective tensor
    with the same positive-length periodic rays as $Q$ has all transfer
    eigenvalues different from $1$ of modulus at most $q$.

    There are a common unit-modulus physical character $\chi$, pointwise
    unital injective tensors $A_t$ of positive bond dimensions bounded by the
    ambient dimension of $Q$, and pointwise unitary projective virtual
    representations with factor systems $\omega_t$, such that $A_t$ has the
    same positive-length periodic rays as $Q_t$, the physical action
    $\chi(g)^{-1}U_g$ has exact inverse-index virtual covariance on $A_t$,
    and $\omega_s$ and $\omega_t$ are cohomologous for all $s,t\in[0,1]$.
    No continuous choice of the tensors, stationary densities, or virtual
    matrices is required. The uniform transfer bound is assumed explicitly;
    this statement does not derive it from the Hamiltonian gap.
\end{lemma}

\begin{proof}\leanok
    \uses{thm:spt_varying_bond_physical_character,
        thm:spc_aux_bounded_minimal_local_class,
        thm:spt_virtual_class_gauge_preparation}
    The finite-chain ground lines determine one common physical character.
    Remove it from the physical action. Normalize each injective periodic-ray
    representative by a nonzero rescaling and bond similarity. Equality of
    periodic rays transfers exact on-site symmetry to these representatives,
    and unitality gives unitary virtual matrices. Injectivity assembles them
    into projective representations with the inverse group convention.
    The bounded-dimension compactness theorem makes the corresponding
    cohomology quotient locally constant. Since the interval is connected,
    this quotient is constant throughout the path.
\end{proof}


\begin{lemma}[A common injective blocking length]
    \label{lem:spc_aux_normal_fourth_power_blocking}
    \lean{MPSTensor.IsNormalTensor.isNBlkInjective_pow_four}\leanok
    As used in arXiv:1606.00608, lines 318--344, and the quantum Wielandt
    bound, for a normalized normal tensor of bond dimension $D$, products of
    length $D^4$ span the full matrix algebra.
\end{lemma}
\begin{proof}\leanok
    Conjugating to a trace-preserving tensor, applying the quantum Wielandt
    bound, and carrying the spanning statement through the bond conjugation
    proves the claim.
\end{proof}

\begin{lemma}[Bounded bond dimension and common physical blocking]
    \label{lem:spc_aux_bounded_normal_fourth_power_blocking}
    \lean{MPSTensor.IsNormalTensor.blockTensor_isInjective_of_bondDim_le}\leanok
    \uses{lem:spc_aux_normal_fourth_power_blocking}
    If a normalized normal tensor has bond dimension $D\le M$, its $M^4$-site
    blocked tensor is injective.
\end{lemma}
\begin{proof}\leanok
    Injectivity at length $D^4$ and preservation of the spanning property
    at every larger length give injectivity at length $M^4$.
\end{proof}

\begin{lemma}[Algebraic simplicity of the fixed eigenvalue]
    \label{lem:spc_aux_positive_unital_simple_fixed_eigenvalue}
    \lean{simple_fixedEigenvalue_of_unital_positive}\leanok
    As the fixed-eigenvalue consequence of Wolf, Proposition 6.2, let $T$ be
    a positive unital linear map on $M_D(\mathbb C)$, with $D>0$: if its fixed
    space is one-dimensional, the eigenvalue $1$ has algebraic multiplicity
    one.
\end{lemma}
\begin{proof}\leanok
    A positive unital map has no nontrivial Jordan block on the unit circle.
    Hence its generalized fixed space is its fixed space, and the algebraic
    multiplicity equals that space's dimension.
\end{proof}

\begin{lemma}[Compactness and class stability for periodic-ray sequences]
    \label{thm:spc_aux_compact_minimal_sequence_class}
    \lean{MPSTensor.eventually_cohomologousTo_of_tendsto_periodic_ray_sequence}
    \leanok
    Let $A$ be an injective unital tensor, and let $B_n$ be unital irreducible
    tensors of one fixed bond dimension. Suppose that their periodic rays are
    represented by tensors $Q_n$ in one fixed ambient dimension, with
    $Q_n\to Q_0$, and that $Q_0$ has the same positive-length periodic rays as
    $A$. Assume $0\le q<1$, and that every transfer eigenvalue different from
    $1$, for $A$ and all $B_n$, has modulus at most $q$. Suppose that one
    physical unitary action has exact inverse-index virtual covariance on
    these tensors, with pointwise unitary projective representatives having
    factor systems $\omega_0$ and $\omega_n$. Then $\omega_n$ is cohomologous
    to $\omega_0$ for all sufficiently large $n$.
\end{lemma}
\begin{proof}\leanok
    \uses{thm:spc_aux_convergent_unital_sequence_class,
        thm:spc_aux_compact_unital_tensors,thm:spc_aux_fixed_line_limit,
        thm:spc_aux_periodic_line_limit}
    If the conclusion fails, choose a subsequence on which every class differs
    from the reference class. Compactness gives a convergent subsequence of
    the unital tensors and their normalized stationary densities. The common
    spectral bound preserves simplicity of the fixed eigenvalue in the limit.
    Nonvanishing and continuity identify the limiting periodic rays with those
    of $A$. The stationary-support sequence theorem gives the contradiction.
\end{proof}

\begin{lemma}[Bounded bond dimensions and local virtual class]
    \label{thm:spc_aux_bounded_minimal_local_class}
    \lean{MPSTensor.eventually_cohomologousTo_of_continuousAt_bounded_unital_periodic_rays}
    \leanok
    Let $T$ be first countable. At each $t\in T$, let $A_t$ be an irreducible
    unital tensor of positive bond dimension $D(t)\le M$. Suppose that fixed
    ambient tensors $Q_t$ represent their positive-length periodic rays and
    vary continuously at $t_0$, and that $A_{t_0}$ is injective. Under one
    exact physical unitary covariance, with pointwise unitary virtual
    representatives, assume a common bound $0\le q<1$ for every nonunit
    transfer eigenvalue. Then the virtual factor systems are cohomologous to
    that at $t_0$ throughout some neighborhood of $t_0$.
\end{lemma}
\begin{proof}\leanok
    \uses{thm:spc_aux_compact_minimal_sequence_class}
    A sequence approaching $t_0$ on which the conclusion fails has a
    subsequence with constant bond dimension; transferring both tensors and
    virtual matrices to that fixed dimension and applying the preceding
    theorem gives the contradiction.
\end{proof}

\begin{lemma}[Positive-length rays under physical blocking]
    \label{lem:spc_aux_positive_ray_blocking}
    \lean{MPSTensor.SamePositiveMpvRay.blockTensor}
    \leanok
    If two tensors have equal periodic rays at every positive length, then
    their $L$-site blocked tensors have equal periodic rays at every positive
    length, for every $L>0$.
\end{lemma}
\begin{proof}\leanok
    Reindexing configurations and applying the original ray equality at
    length $NL$ yields the blocked ray equality.
\end{proof}

\begin{lemma}[Normal realizations of bounded dimension]
    \label{thm:spc_aux_bounded_normal_local_class}
    \lean{MPSTensor.eventually_cohomologousTo_of_continuousAt_bounded_normal_periodic_rays}
    \leanok
    Let $T$ be first countable, and let normal unital tensors $A_t$ have
    bounded bond dimensions $D(t)\le M$. Suppose that fixed ambient tensors
    $Q_t$ represent their positive-length periodic rays and vary continuously
    at $t_0$. Assume exact covariance under a common physical unitary action,
    pointwise unitary virtual representatives, and a common bound
    $0\le q<1$ on all transfer eigenvalues different from $1$. Then their
    virtual factor systems are locally cohomologous to that at $t_0$.
    Neither minimal tensors nor stationary densities need be chosen
    continuously. The uniform transfer bound is an explicit hypothesis.
\end{lemma}
\begin{proof}\leanok
    \uses{thm:spc_aux_bounded_minimal_local_class,
        lem:spc_aux_bounded_normal_fourth_power_blocking,
        lem:spc_aux_positive_ray_blocking,thm:spc_aux_blocked_unitary_covariance}
    Blocking $M^4$ sites, normality and the quantum Wielandt bound imply
    injectivity at every parameter. Blocking preserves the virtual factor
    system and unitality, and changes the spectral bound to $q^{M^4}<1$, so
    the bounded-dimension local theorem applies.
\end{proof}


% Auxiliary preparation for compact-sequence class stability.
% The transfer contraction bound is supplied throughout.

\begin{lemma}[Uniqueness of a normalized adjoint stationary matrix]
    \label{thm:spc_aux_normalized_stationary_unique}
    \lean{MPSTensor.normalized_adjoint_fixed_unique_of_transfer_fixed_finrank_one}
    \leanok
    \uses{def:transfer_map}
    As an auxiliary identity for
    \cite[Appendix~C, lines~2653--2717]{Schuch2011Phases}, let $B$ have
    arbitrary finite bond dimension and $\dim_{\C}\ker(\E_B-\Id)=1$.
    If $\E_B^*(\sigma)=\sigma$, $\E_B^*(Z)=Z$, and
    $\tr\sigma=\tr Z=1$, then $Z=\sigma$, without requiring positivity or
    unitality.
\end{lemma}
\begin{proof}\leanok
    \uses{thm:spc_aux_adjoint_fixed_dimension}
    Trace duality makes the adjoint fixed space one-dimensional.
    Since $\tr\sigma=1$, the matrix $\sigma$ is nonzero and spans that space.
    Thus $Z=c\sigma$; taking traces gives $c=1$.
\end{proof}

\begin{lemma}[Virtual class under gauge equivalence up to a scalar]
    \label{thm:spc_aux_injective_gauge_class}
    \lean{MPSTensor.cohomologousTo_of_injective_gaugePhaseEquiv}
    \leanok
    \uses{def:gauge_phase_equiv, def:projective_rep,
        def:cohomologous_to}
    Following \cite[Section~II.F.2 and Appendix~C]{Schuch2011Phases}, let
    $A,B$ have the same positive bond dimension, with $A$ injective,
    and suppose $B^i=cXA^iX^{-1}$ for $c\ne0$ and $X$ invertible.
    Let a physical group representation $U$ have virtual projective
    representations $V_A,V_B$ with factors $\omega_A,\omega_B$ satisfying
    $\sum_jU(g)_{ij}A^j=V_A(g^{-1})A^iV_A(g^{-1})^{-1}$ and the
    corresponding identity for $B$. Then $\omega_B\sim\omega_A$.
\end{lemma}
\begin{proof}\leanok
    \uses{thm:spt_virtual_class_gauge_preparation}
    The scalar and gauge preserve injectivity. Conjugating $V_A$ by $X$
    retains its factor and implements the symmetry of $B$.
    Gauge independence on the injective tensor $B$ compares that action with
    $V_B$ and gives the asserted cohomology relation.
\end{proof}

\begin{lemma}[Unitary virtual transfer between unital tensors]
    \label{thm:spc_aux_unitary_gauge_virtual_transport}
    \lean{MPSTensor.exists_unitary_virtualRep_of_unital_gaugePhaseEquiv}
    \leanok
    \uses{def:gauge_phase_equiv, def:projective_rep}
    Following \cite[Appendix~C, lines~2653--2717]{Schuch2011Phases}, let
    $A,B$ have the same positive bond dimension, be unital, and be
    gauge equivalent up to a nonzero scalar, with $A$ injective.
    Suppose a unitary physical group representation $U$ has a unitary
    projective virtual action $V$ on $A$, with factor $\omega$, satisfying
    $\sum_jU(g)_{ij}A^j=V(g^{-1})A^iV(g^{-1})^\dagger$.
    Then $B$ admits a unitary projective virtual action with exactly the
    factor $\omega$ and the same covariance identity.
\end{lemma}
\begin{proof}\leanok
    \uses{lem:pgvwc07_intertwiner_unitary, thm:spt_virtual_class_gauge_preparation}
    The scalar and gauge make $B$ injective. The unital intertwiner theorem
    replaces the bond gauge by a unitary $W$, giving
    $B^i=cWA^iW^\dagger$. The representation $WV(g)W^\dagger$ is unitary,
    retains the factor $\omega$, and satisfies the required covariance
    under the scalar and bond gauge.
\end{proof}

\begin{definition}[Unitary on-site action after blocking]
    \label{def:spc_aux_block_unitary_action}
    \lean{MPSTensor.blockUnitaryPhysicalAction}
    \leanok
    \uses{def:block_kron_action}
    Let $U$ be a monoid representation by unitary matrices on $\C^d$.
    For $L\geq0$, its unitary action on the blocked physical space is
    $U^{[L]}(g)=U(g)^{\otimes L}$, in the coordinates of length-$L$ words.
\end{definition}
\begin{proof}\leanok
    \uses{lem:on_site_tensor_pow_mul}
    Tensor powers preserve multiplication and identity. They commute with
    adjoints, so $U^{[L]}(g)U^{[L]}(g)^\dagger=\Id$.
\end{proof}

\begin{lemma}[Unchanged virtual factor under common blocking]
    \label{thm:spc_aux_blocked_unitary_covariance}
    \lean{MPSTensor.rotatePhysical_blockTensor_of_unitary_virtual_covariance}
    \leanok
    \uses{def:spc_aux_block_unitary_action, def:blocked_tensor, def:projective_rep}
    Following \cite[Appendix~C, lines~2653--2717]{Schuch2011Phases}, suppose
    a unitary physical group representation $U$ and a unitary
    projective virtual representation $V$, with factor $\omega$, satisfy
    $\sum_jU(g)_{ij}A^j=V(g^{-1})A^iV(g^{-1})^\dagger$.
    For every $L\geq0$, the same $V$ implements the covariance of $A^{[L]}$
    under $U^{[L]}$. In particular, its factor remains exactly $\omega$.
\end{lemma}
\begin{proof}\leanok
    \uses{lem:twisted_tensor_block_gauge}
    Twisting commutes with blocking. Word evaluation of the virtual gauge
    telescopes, leaving $V(g^{-1})A^wV(g^{-1})^{-1}$.
    Unitarity replaces the last inverse by the adjoint. The representation
    itself has not changed, so neither has its factor.
\end{proof}

\begin{lemma}[Overlap modulus of equal periodic spans]
    \label{thm:spc_aux_periodic_span_overlap}
    \lean{MPSTensor.norm_mpvOverlap_eq_mul_norm_of_span_eq}
    \leanok
    \uses{def:mpv_state, def:mpv_overlap}
    Let $A,B$ have the same physical dimension and arbitrary bond dimensions.
    At any $N\geq0$, equality of the spans of their periodic coefficient
    vectors implies
    $|O_{AB}(N)|=\|\ket{V^{(N)}(A)}\|\,\|\ket{V^{(N)}(B)}\|$.
    Both vectors are allowed to vanish.
\end{lemma}
\begin{proof}\leanok
    \uses{lem:overlap_eq_inner}
    Equality of spans gives $\ket{V^{(N)}(B)}=c\ket{V^{(N)}(A)}$.
    Substitution into the inner product gives
    $|O_{AB}(N)|=|c|\|\ket{V^{(N)}(A)}\|^2$, which is the asserted
    product by homogeneity of the norm.
\end{proof}

\begin{lemma}[Unit limiting periodic norm from transfer contraction]
    \label{thm:spc_aux_periodic_norm_one}
    \lean{MPSTensor.norm_mpvState_tendsto_one_of_transfer_spectrum}
    \leanok
    \uses{def:transfer_map, def:mpv_state}
    Let $A$ be unital with positive bond dimension $D$ and
    $\dim_{\C}\ker(\E_A-\Id)=1$.
    Suppose $0\leq q<1$ and $|z|\leq q$ for every
    $z\in\operatorname{spec}(\E_A)$ with $z\ne1$.
    Then $\|\ket{V^{(N)}(A)}\|\to1$.
    The spectral bound is an explicit auxiliary assumption in the setting
    of \cite[Appendix~C, lines~2653--2717]{Schuch2011Phases}.
\end{lemma}
\begin{proof}\leanok
    \uses{thm:spt_compact_periodic_norm_deviation}
    The periodic norm estimate gives
    $|\|\ket{V^{(N)}(A)}\|^2-1|\leq(D^2-1)q^N$.
    Since $q^N\to0$, the squared norms tend to one. Continuity of the square
    root and nonnegativity of the norms finish the proof.
\end{proof}

\begin{lemma}[Unit overlap limit from eventual periodic-line equality]
    \label{thm:spc_aux_periodic_line_unit_overlap}
    \lean{MPSTensor.norm_mpvOverlap_tendsto_one_of_eventual_span_eq}
    \leanok
    \uses{def:mpv_state, def:mpv_overlap}
    Let $A,B$ have the same physical dimension and arbitrary bond dimensions.
    Suppose their periodic coefficient spans agree for all sufficiently
    large $N$, and both periodic state norms tend to one.
    Then $|O_{AB}(N)|\to1$.
\end{lemma}
\begin{proof}\leanok
    \uses{thm:spc_aux_periodic_span_overlap}
    Eventually the overlap modulus is the product of the two state norms.
    That product tends to one. No proportionality scalars are chosen
    simultaneously across the chain lengths.
\end{proof}


\begin{lemma}[Virtual class along a convergent unital sequence]
    \label{thm:spc_aux_convergent_unital_sequence_class}
    \lean{MPSTensor.eventually_cohomologousTo_of_tendsto_unital_periodic_lines}
    \leanok
    \uses{def:projective_rep, def:transfer_map, def:mpv_state,
        def:cohomologous_to}
    Let $A$ be a unital injective tensor with positive bond dimension $D$.
    Let $B_n$ be unital tensors of a fixed bond dimension $k$, converging to
    a unital tensor $b$ with $\dim_{\C}\ker(\E_b-\Id)=1$.
    Suppose $B_n$ has a positive-definite adjoint stationary matrix $\sigma_n$
    of trace one. Let $\tau$ be a positive semidefinite adjoint stationary
    matrix for $b$, with trace one, and let $K:\C^E\to\C^k$ be an isometry
    whose range is the support of $\tau$.
    Write $C^i=K^\dagger b^iK$.

    Suppose $0\leq q<1$ bounds the modulus of every nonunit eigenvalue of
    both $\E_A$ and $\E_b$. Suppose also that the periodic coefficient spans
    of $A$ and $b$ agree for all sufficiently large lengths.
    Let the same unitary physical group representation $U$ act covariantly
    on $A$ and every $B_n$, through pointwise unitary projective virtual
    representations with respective factors $\omega_0$ and $\omega_n$:
    $\sum_j U(g)_{ij}A^j=V_0(g^{-1})A^iV_0(g^{-1})^\dagger$, and likewise
    for $B_n$ and $V_n$.
    Then $\omega_n\sim\omega_0$ for all sufficiently large $n$.
    No continuity of the virtual representations or of $\sigma_n$ is required.
    The transfer bound and periodic-span equality are explicit hypotheses;
    this is an auxiliary reduction for
    \cite[Appendix~C, lines~2653--2717]{Schuch2011Phases}.
\end{lemma}
\begin{proof}\leanok
    \uses{thm:injective_irreducible,
        thm:spt_unitary_virtual_covariance_unital,
        thm:spc_aux_periodic_norm_one,
        thm:spc_aux_periodic_line_unit_overlap,
        thm:spc_aux_stationary_support_identification,
        thm:spc_aux_normalized_stationary_unique,
        thm:spc_aux_stationary_support_preparation,
        thm:spc_aux_unitary_gauge_virtual_transport,
        thm:spc_aux_adjoint_fixed_dimension,
        thm:spc_aux_topological_stationary_support_class}
    Injectivity makes $\E_A$ irreducible. Since $\E_A(I)=I$,
    Perron--Frobenius theory makes its fixed space one-dimensional,
    and unital normalization gives a positive definite matrix
    $\Lambda$ with $\E_A^\dagger(\Lambda)=\Lambda$.
    The two periodic norms tend to one. Equality of their eventual spans
    therefore gives a unit overlap limit. Stationary-support identification,
    applied with the irreducible family $A$ and the fixed point $\Lambda$,
    shows that $E=D$ and that $C$ is gauge equivalent to $A$ up to a nonzero
    scalar. In particular, $C$ is injective. The fixed-space hypothesis gives
    uniqueness of the normalized adjoint stationary matrix, and support
    preparation makes $C$ unital. Unitary gauge transfer gives a virtual
    action on $C$ with exactly the factor $\omega_0$.

    Adjoin the limiting point to the sequence. The ambient tensors form a
    continuous family on this one-point compactification. At the finite points,
    their stationary support is the full bond space; at the limiting point
    it is the range of $K$. The local stationary-support theorem applies
    without continuous support frames or continuous virtual representatives.
    Its conclusion at the limiting point is precisely the asserted eventual
    cohomology relation.
\end{proof}


\begin{lemma}[Endpoint cohomology for a given unital family]
    \label{thm:spt_unital_supported_path_class}
    \lean{TNLean.Algebra.isMulCoboundary₂_div_of_eventually_at_all_points,
        MPSTensor.H2π_eq_of_continuous_unital_supported_family}
    \leanok
    In the unital case of the preceding local theorem, suppose the ambient
    adjoint fixed space has complex dimension one and every minimal support
    compression is one-site injective at every parameter in $[0,1]$.
    Then the endpoint virtual factors are cohomologous. More generally,
    a scalar cohomology class that is locally constant on a preconnected
    parameter space is constant throughout that space.
\end{lemma}

\begin{proof}\leanok
    \uses{thm:spt_local_unique_stationary_class}
    Local constancy gives a locally constant map into the quotient by
    scalar cohomology. Such a map is constant on a preconnected space.
    Apply this to the unit interval and its two endpoints.
\end{proof}
